\documentclass[11pt,a4paper,reqno]{amsart}
\PassOptionsToPackage{expansion=false}{microtype}
\usepackage[T1]{fontenc}
\usepackage{lmodern,mathtools,amssymb,amsthm,mathrsfs,booktabs,microtype}
\usepackage[a4paper,textwidth=156mm,textheight=237mm,centering]{geometry}
\usepackage[hidelinks]{hyperref}
\hypersetup{pdftitle={Rational approximation to power-sum quotients and simultaneous recurrence},pdfauthor={},pdfkeywords={Subspace Theorem, power sums, rational approximation, linear recurrences, Hausdorff dimension}}
\newtheorem{theorem}{Theorem}[section]
\newtheorem{lemma}[theorem]{Lemma}
\newtheorem{proposition}[theorem]{Proposition}
\newtheorem{corollary}[theorem]{Corollary}
\theoremstyle{definition}
\newtheorem{definition}[theorem]{Definition}
\newtheorem{example}[theorem]{Example}
\theoremstyle{remark}
\newtheorem{remark}[theorem]{Remark}
\numberwithin{equation}{section}
\setcounter{tocdepth}{1}
\setlength{\emergencystretch}{1.4em}
\title[Power-sum quotients and simultaneous recurrence]{Rational approximation to power-sum quotients\\and simultaneous recurrence}
\date{}
\raggedbottom
\begin{document}
\begin{abstract}
We prove Roth's exponent for quotients of positive integer-root power sums under the hypothesis of Corvaja and Zannier. Every fixed denominator range \(q\le Q^n\) with \(1<Q<2\) is admissible, and \(2\) is the sharp universal boundary. For vectors of \(m\) quotients we characterize the sharp range at exponent \(1+1/m\) by rational-function subspaces. The proof computes the weight of a box of monomial multiples after restriction to a subspace of arbitrary rank. As an application we determine the local Hausdorff dimension of simultaneous recurrence on the diagonal for any number of pairwise coprime bases, separate shrinking radii and level-dependent translations. We also give a scalar cutoff for algebraic linear recurrences at a prescribed embedding, and a cutoff for approximation by polynomials of bounded degree. For the trace ratio of the cubic unit of minimal polynomial \(X^3-X-1\), its real root is the sharp height cutoff for quadratic approximation.
\end{abstract}
\maketitle
\tableofcontents

\section{Introduction}
Let \(\mathcal E\) be the ring of sequences \(u(n)=\sum_{j=1}^t c_j a_j^n\), with \(c_j\in\mathbb Q\) and positive integers \(a_j\). After combining equal roots and discarding zero coefficients, write \(\ell(u)\) for the largest root and put \(\ell(0)=0\). Corvaja and Zannier \cite[pp.~738--740]{CZ05} proved a Thue-type inequality for the moving rational number \(u(n)/v(n)\). Their denominator exponent and the base of the exponential denominator range may both depend on \(u,v\). In Remark~(v) on p.~740 they asked for Roth's exponent. We obtain it under their original hypothesis.
\begin{theorem}\label{thm:cz}
Let \(u,v\in\mathcal E\) be nonzero. Suppose that
\begin{equation}\label{eq:czhyp}
                 \ell(u-\zeta v)\ge\ell(v)\qquad(\zeta\in\mathcal E).
\end{equation}
For every fixed \(1<Q<2\) and every \(\epsilon>0\), all sufficiently large \(n\) satisfy
\[
 \left|\frac{u(n)}{v(n)}-\frac pq\right|\ \ge\ \frac{e^{-\epsilon n}}{q^2}
 \qquad(p,q\in\mathbb Z,\ 1\le q\le Q^n).
\]
The exponent \(2\) is optimal for every such \(Q\), and \(Q=2\) can fail.
\end{theorem}
The cutoff for an individual pair can be much larger. Theorem~\ref{thm:scalarcut} computes it by finitely many linear systems over \(\mathbb Q\) and characterizes equality using reduced denominators. In all our cutoff statements the upper bound on height or denominator includes its endpoint unless specified otherwise. Exact approximants count as violations. A conclusion holding eventually is uniform in the approximants for the fixed input sequences and tolerance.

The arithmetic input is Theorem~\ref{thm:descent}. Take the primes in the integer roots as generators, write the orbit as \(X_n=(a_1^n,\ldots,a_d^n)\), and regard the moving targets as rational functions of \(X_n\). Multiplying numerator and denominator coordinates of an approximation by a box of monomials gives a system to which the ordinary \(S\)-adic Subspace Theorem applies. For several targets, a first relation can leave a subspace of dimension greater than one. We compute the weight of the remaining forms on that subspace in terms of its Pl\"ucker coordinates. This gives a descent at each rank and identifies the subspaces whose rescaled lattice covolumes permit exceptional approximations. The resulting criterion at exponent \(1+1/m\) appears in Theorem~\ref{thm:simcut}.

For instance, set \(R_n=((3^n-1)/(2^n-1),(5^n-1)/(2^n-1))\). The estimate \(q^{1/2}\max_i\|qR_{i,n}\|\ge e^{-\epsilon n}\) holds below the exact common denominator. Its sharp exponential cutoff is \(2\). For the dual quantity \(\|h_1R_{1,n}+h_2R_{2,n}\|\) the corresponding coefficient-height cutoff is \(\sqrt2\) (Theorems~\ref{thm:primal} and~\ref{thm:dual}). Here \(\|x\|\) is distance to the nearest integer. These estimates govern returns of one parameter under several multiplying maps: \(b_i^nx\) returns close to \(x+y_{i,n}\) on the circle precisely when \((b_i^n-1)x\) comes close to \(y_{i,n}\). For ordered pairwise coprime bases and a common radius \(e^{-gn}\), the local dimension is
\[
              \left[\min_{1\le j\le k}
              \frac{\log b_j-(j-1)g}{\log b_j+g}\right]_+ ,
              \qquad [z]_+=\max(z,0).
\]
Theorem~\ref{thm:rec} allows separate nonmonotone radii and all moving phases. Common prime factors can change both the dimension and its dependence on the phases.

Two extensions use different heights. The scalar number-field theorem (Theorem~\ref{thm:algcut}) concerns the rational denominator \(q\) at a chosen embedding, including targets formed from polynomial multiples of algebraic powers. Theorems~\ref{thm:polycut}, \ref{thm:ratpower} and~\ref{thm:cubicquad} concern the maximum absolute coefficient of an integral polynomial or of a primitive minimal polynomial. At \(R_n=(3/2)^n\), the exact rational approximant has denominator \(2^n\) and minimal-polynomial height \(3^n\). The two boundaries in this example are therefore distinct even in degree one.

\subsection*{Relation to previous work}
Corvaja and Zannier \cite[p.~738]{CZ05} recall the exponent-two result for a single rational power. Their Corollary~3 proves divergence of the continued-fraction length for \((b^n-1)/(a^n-1)\) when \(a,b\) are multiplicatively independent. Their later book returns to this example \cite[Exercise~4.47, pp.~165--166]{CZ18}. Fuchs and Heintze \cite[Theorems~2.1 and~2.6]{FH} treat rational approximation to values of algebraic functions over power sums. Their exponent depends on the data, with the denominator range linked to the tolerance. In Theorem~\ref{thm:cz} the same \(Q\) works for every positive tolerance. Kumar and Thangadurai \cite[Theorem~1.1]{KT} treat a vector of fixed algebraic multiples of a common moving group element. The coordinates in Theorem~\ref{thm:simcut} can have different numerator functions.

The outside-\(S\) gcd estimate used below is due to Corvaja and Zannier in two variables \cite{CZgcd}, and to Levin in higher dimensions \cite{Levin}; see the arithmetic formulation in \cite[Theorem~1.13(a)]{LW}. Xiao's estimate \cite{Xiao} also accommodates polynomial factors in the index. The method of passing from moving forms to fixed forms in a larger space has antecedents in the moving-target Subspace Theorem \cite{RV}. Here the coefficient heights may be of the same order as the heights of the approximants.

For a fixed real algebraic target of degree greater than \(d\), Schmidt's theorem yields the bounded-degree approximation exponent \(d+1\), up to an arbitrarily small loss in the power of the height \cite{Schmidt70}. The polynomial-value and root-distance exponents for fixed targets are studied in \cite{Wirsing,BE}. In Section~\ref{sec:higher-degree} the target changes with \(n\), the loss is \(e^{-\epsilon n}\), and the height range is part of the conclusion.

There are two limitations on the bibliographical comparison. We had partial access to the 2008 survey of Corvaja and Zannier \cite{CZ08}. Kumar and Prasad Sena list a separate submitted manuscript \cite{KPsub}; the public research summary \cite[p.~59]{HRI} describes nearest-integer and continued-fraction applications for quotients, but we could not verify its general theorem. Their published result on linear combinations is \cite[Theorem~1.3]{KP}. We leave priority with respect to the unavailable statements open.



\section{Descent on a power orbit}\label{sec:descent}

Let \(d\geq 1\), let \(a_1,\ldots,a_d\geq 2\) be pairwise coprime
integers, and put \(X_n=(a_1^n,\ldots,a_d^n)\),
\(\mathcal P=\mathbb Q[X_1,\ldots,X_d]\), and
\(F=\operatorname{Frac}(\mathcal P)\). For \(f\in\mathcal P\) put
\[
 w(f)=\max_{\alpha\in\operatorname{supp}(f)}
           \sum_{j=1}^d \alpha_j\log a_j,\qquad w(0)=-\infty .
\]
The weights of distinct monomials are distinct by pairwise coprimality.
In particular \(w(fg)=w(f)+w(g)\), so \(w\) extends to \(F\) by
\(w(f/g)=w(f)-w(g)\) for \(g\ne 0\). For nonzero \(f\in F\),
\begin{equation}\label{eq:orbit-weight}
 \log |f(X_n)|=nw(f)+O_f(1)
\end{equation}
at all sufficiently large \(n\). If \(A\subset F^k\) is a subspace,
its primitive polynomial Pl\"ucker vector is a generator of
\(\bigwedge^{\dim A} A\) with polynomial entries of polynomial gcd \(1\).
It exists by clearing denominators and dividing out the gcd, and is
unique up to \(\mathbb Q^\times\). In dimension zero we use the entry
\(P_{\varnothing}=1\). For a fixed \(F\)-subspace \(W\), we specialize
a polynomial basis of the equations defining \(W\) and write \(W(X_n)\)
for the common kernel. The resulting subspace is independent of the
basis and has dimension \(\dim_F W\) for all sufficiently large \(n\),
by \eqref{eq:orbit-weight} applied to the relevant denominators and minors.

\begin{theorem}[Descent]\label{thm:descent}
Let \(W\subset F^k\) have dimension \(r\geq 1\), put \(c=k-r\), and let
\(P\) be the primitive polynomial Pl\"ucker vector of \(W^\perp\).
Let \(S\in\mathrm{GL}_k(F)\) keep a subset of the coordinates \(x_j\) and
replace each remaining coordinate by
\[
 (Sx)_i=x_i-\sum_{\text{\(j\) kept}}\frac{U_{ij}}{V_i}x_j,
 \qquad U_{ij},V_i\in\mathcal P,\qquad V_i\ne0 .
\]
The identity is allowed. Suppose that, on an infinite set of levels
\(n\), there are nonzero integer vectors \(x_n\in W(X_n)\) with
\(\max_i |x_{n,i}|\leq e^{Cn}\) for some fixed \(C\), and
\[
 |(S(X_n)x_n)_i|\leq \exp\bigl(n(\mu_i+o(1))\bigr)
\]
for fixed real numbers \(\mu_i\). Define
\begin{equation}\label{eq:descent-K}
 \mathcal B=\sum_{i=1}^k\mu_i,\qquad
 \mathcal K=\max_{|I|=c}
    \left\{\sum_{i\in I}\mu_i+
           w\bigl((\bigwedge^c S^{-T}P)_I\bigr)\right\}.
\end{equation}
If \(\mathcal K>\mathcal B\), a fixed proper \(F\)-subspace of \(W\)
contains \(x_n\) upon specialization for infinitely many of these levels.
Equivalently, writing \(P^W\) for the primitive polynomial Pl\"ucker
vector of \(W\), the strict inequality holds if and only if
\begin{equation}\label{eq:descent-primal}
 \max_{|J|=r}
   \left\{w\bigl((\bigwedge^r S P^W)_J\bigr)
           -\sum_{i\in J}\mu_i\right\}>0.
\end{equation}
The maximum in \eqref{eq:descent-primal} equals \(\mathcal K-\mathcal B\).
\end{theorem}

The equality of the two expressions follows from complementary
minors: the entries of \(P^W\) are, up to signs and a common rational
scalar, the complementary entries of \(P\), and \(\det S=1\).
We prove descent by applying the fixed-form \(S\)-adic Subspace
Theorem to a finite-dimensional space of polynomial multiples of
the coordinate forms. The following lemma computes the total weight
of the relations among these forms.

\subsection{Weights of polynomial relations}

For a real vector \(\lambda=(\lambda_1,\ldots,\lambda_d)\) and coordinate
costs \(\nu=(\nu_1,\ldots,\nu_k)\), give the monomial \(X^\alpha e_i\)
weight \(\lambda\cdot\alpha+\nu_i\), where \(e_i\) is the \(i\)-th
coordinate vector. A polynomial vector has the largest weight of its
nonzero terms, with weight \(-\infty\) for zero. On a subspace \(L\) of
a finite-dimensional polynomial space, its pivot weights are the
jumps, counted with multiplicity, of the dimension of the filtration
by weight at most \(t\). Write \(\Omega_{\lambda,\nu}(L)\) for their sum.

For a quotient, give each nonzero class the smallest weight of a
representative. For a filtered space and any subspace, the multiset
of weights of the space is the union of those of the subspace and
quotient, by additivity of dimension at each \(t\). We use the same
definitions for weights specified by other bases. If two weight
functions on a space differ by at most \(K\) on nonzero vectors,
their ordered pivot weights differ termwise by at most \(K\), since
each filtration is contained in the other filtration shifted by \(K\).
The same comparison holds for the induced quotient weights.

\begin{lemma}[Box filtration]\label{lem:box-filtration}
Let \(A\subset F^k\) have dimension \(c\), with primitive polynomial
Pl\"ucker vector \(P\). Let \(\mathcal P_h\) be the polynomials supported
in \([0,h]^d\) and put
\[
 A_h=A\cap\mathcal P_h^k,\qquad D_h=\dim_{\mathbb Q} A_h,\qquad
 M_h=(h+1)^d,\qquad \overline w_h=\frac h2\sum_{j=1}^d\lambda_j .
\]
Write \(d_j=\max_{I:P_I\ne0}\deg_{X_j}P_I\) and set
\[
 K_{\lambda,\nu}(P)=
 \max_{\substack{|I|=c\\P_I\ne0}}
   \left\{\sum_{i\in I}\nu_i+
          \max_{\alpha\in\operatorname{supp}(P_I)}\lambda\cdot\alpha\right\}.
\]
Then \(D_h=cM_h+O(h^{d-1})\), and
\begin{equation}\label{eq:box-filtration}
 \Omega_{\lambda,\nu}(A_h)
   =D_h\overline w_h+
     M_h\left(K_{\lambda,\nu}(P)-\frac12\sum_{j=1}^d\lambda_jd_j\right)
     +O(h^{d-1}).
\end{equation}
\end{lemma}

\begin{proof}
The assertion is immediate for \(c=0\), so suppose \(c\geq1\).
Use the multigraded ring in \(2d\) variables \(T_{j,0},T_{j,1}\),
each of degree \(e_j\) in its pair, where \(e_j\) here denotes the
\(j\)-th standard multidegree. In degree \((h_1,\ldots,h_d)\)
identify the monomial
\(\prod_j T_{j,0}^{h_j-\alpha_j}T_{j,1}^{\alpha_j}\) with \(X^\alpha\).
The homogeneous lifts of the polynomial vectors of \(A\) in all
boxes form a multigraded submodule of the free module of rank \(k\).
Indeed, \(T_{j,0}\) and \(T_{j,1}\) act after dehomogenization by
multiplication by \(1\) and \(X_j\), respectively.

Order terms first by their weights and then lexicographically by
the exponent vector in the \(T\)'s and the coordinate tag vector
\(e_i\in\mathbb Z^k\). On wedge terms use the sum of the tags of the
participating coordinates and the sum of their costs. This comparison
is additive, and strict inequalities persist when the same data
are added to both sides. In particular, if vectors have distinct
leading coordinates then their wedge has as leading term the
product of their leading terms, with the corresponding wedge tag.
That product is nonzero and all other nonzero products are strictly
smaller. These comparisons are used on finite supports, so no
well-ordering assumption on the weights is needed.

The leading terms of homogeneous elements generate a monomial
submodule \(\bigoplus_i J_i e_i\). Each of its monomials is itself
a leading monomial, by compatibility with monomial multiplication.
In any fixed multidegree, row reduction in the ordered monomial
basis shows that these leading monomials give precisely the pivot
slots, with their filtration weights. There are at most \(c\)
nonzero \(J_i\), since vectors with distinct leading coordinates
have a nonzero wedge and are independent over \(F\) after
dehomogenizing. Conversely, take \(c\) independent polynomial
vectors of \(A\). Their monomial multiples with exponents in
\([0,h-M]^d\) lie in \(A_h\) for a fixed \(M\), and are
\(\mathbb Q\)-linearly independent. Each component contributes at
most \(M_h\) pivots, so there must be \(c\) nonzero \(J_i\).

For each nonzero \(J_i\), let
\[
 M_i=\prod_{j,\sigma}T_{j,\sigma}^{e_{i,j,\sigma}}
\]
be the monomial gcd of its monomials. After dividing the ideal by
\(M_i\), for every variable there is a monomial generator with
exponent zero in that variable. Fix a finite set of generators,
and fix a bound on all their exponents. A monomial whose exponents
are at least this bound except possibly in one variable is divisible
by a generator. Thus the monomials of the principal ideal \((M_i)\)
missing from \(J_i\) lie within fixed distances of at least two
distinct faces of the exponent box. With all lengths sufficiently
large, these faces must be on different axes, since the exponents
of \(T_{j,0}\) and \(T_{j,1}\) sum to \(h_j\). The number missing is
\(O(h^{d-2})\) when all lengths are \(h\) and \(d\geq 2\). When
\(d=1\) there are no missing monomials for large \(h\). In particular
the diagonal count is \(D_h=cM_h+O(h^{d-1})\).

We next show that
\begin{equation}\label{eq:box-face-sum}
 \sum_{i:J_i\ne0}(e_{i,j,0}+e_{i,j,1})=d_j
 \qquad (1\leq j\leq d).
\end{equation}
Fix \(h_j=H\) sufficiently large and set all other lengths to \(h\).
The coefficient of \(h^{d-1}\) in the pivot count is
\(c(H+1)-\sum_{i:J_i\ne0}(e_{i,j,0}+e_{i,j,1})\)
in the sense that the count equals this coefficient times \(h^{d-1}\)
plus \(O_H(h^{d-2})\) for \(d\ge2\). Indeed the shifted box in component
\(i\) has \(H+1-e_{i,j,0}-e_{i,j,1}\) slots on axis \(j\),
and the two-face omissions restrict at least one other axis.
For \(d=1\) the count at large \(H\) is exactly the displayed coefficient.

Let \(K_j=\mathbb Q(X_\ell:\ell\ne j)\). The same coefficient is the
dimension over \(K_j\) of the vectors of \(A\) in \(K_j[X_j]^k\) with
\(X_j\)-degree at most \(H\). To see the lower bound, clear the fixed
denominators in a \(K_j\)-basis of these vectors and multiply them
by monomials in the other variables of degrees at most \(h-O_H(1)\).
For the upper bound, project injectively on this \(K_j\)-space onto
as many coordinate coefficient slots in \(X_j\) as its dimension.
For polynomial vectors in the given box, each such slot is a
polynomial in the other variables supported in \([0,h]^{d-1}\).
The two bounds give the claimed coefficient. For \(d=1\) this
dimension is the original count over \(\mathbb Q\).

The module \(A\cap K_j[X_j]^k\) is saturated, since \(A\) is a vector
space over \(K_j(X_j)\). Over the principal ideal domain \(K_j[X_j]\),
it is free of rank \(c\), and its quotient in the ambient module is
free, being finitely generated and torsion-free. Consequently its
basis extends to an ambient basis and the wedge of its basis vectors
is primitive. Choose such a basis with linearly independent leading
coefficient vectors. To obtain it, in any dependence of the leading
coefficient vectors choose a participating row of maximal degree.
Add suitably shifted multiples of the other participating rows to
that row to cancel its leading coefficient. This is a unimodular
operation and strictly decreases the sum of row degrees, so it
terminates. For row degrees \(\delta_i\), the degree of a nonzero
combination with polynomial coefficients \(f_i\) is then
\(\max_i(\delta_i+\deg f_i)\). For large \(H\) the dimension is
\(c(H+1)-\sum_i\delta_i\). The degree of the basis wedge is
\(\sum_i\delta_i\) because the leading coefficient wedge is nonzero.
The original \(P\) remains primitive over \(K_j[X_j]\), as localization
of the polynomial unique factorization domain can only remove
irreducible common factors. The two primitive generators therefore
differ by a unit in \(K_j[X_j]\), giving \(\sum_i\delta_i=d_j\).
This proves \eqref{eq:box-face-sum}.

Multihomogenize \(P\) to degree \((d_1,\ldots,d_d)\).
The resulting tuple is primitive in the \(T\)-ring. A common
irreducible factor of its homogeneous entries is multihomogeneous
and, unless it is one of the \(T_{j,0}\), dehomogenizes to a common
nonconstant factor of \(P\). None of the \(T_{j,0}\) is a common
factor by the definition of the respective \(d_j\).
For each fixed variable \(T_{j,\sigma}\), choose one homogeneous
form in every nonzero leading component that attains the minimum
exponent \(e_{i,j,\sigma}\) in its leading term. The wedge of these
forms is a polynomial multiple of the multihomogenized \(P\):
dehomogenization gives proportionality, multihomogeneity gives
proportionality over the fraction field of the \(T\)-ring, and
primitivity clears any denominator. The leading term of the wedge
is the product of the chosen leading terms. It is also the product
of the leading term of \(P\) and that of the scalar multiple, by
additivity of the comparison. Hence the \(T_{j,\sigma}\)-exponent
in the leading term of \(P\) is at most \(\sum_i e_{i,j,\sigma}\).
Since the exponent sum for each pair is \(d_j\), the two bounds are
equalities by \eqref{eq:box-face-sum}. The same comparison identifies
the index of the leading term of \(P\) with the set of nonzero leading
components.

In diagonal degree \(h\), the mean weight in the shifted box of
component \(i\) is
\[
 \overline w_h+\nu_i+
   \frac12\sum_j\lambda_j(e_{i,j,1}-e_{i,j,0}).
\]
Subtract \(\overline w_h\) from the weight of each pivot. The missing
monomials give an error \(O(h^{d-1})\), since there are
\(O(h^{d-2})\) of them of centered weight \(O(h)\), and none when
\(d=1\) and \(h\) is large. Each shifted box has
\(M_h+O(h^{d-1})\) slots, and its centered mean is constant in \(h\).
The exponent and coordinate identities just proved show that the
sum of these centered means is
\(K_{\lambda,\nu}(P)-\frac12\sum_j\lambda_jd_j\). Adding back
\(D_h\overline w_h\) gives \eqref{eq:box-filtration}.
\end{proof}


\subsection{Local bases and proof of descent}

\begin{proof}[Proof of Theorem~\ref{thm:descent}]
Let \(\mathcal S_f\) be the set of primes dividing \(\prod_j a_j\)
and \(\mathcal S=\mathcal S_f\cup\{\infty\}\). Normalize the finite
absolute values by \(|p|_p=p^{-1}\), and put
\(\lambda_{v,j}=\log|a_j|_v\). Then
\(\sum_{v\in\mathcal S}\lambda_{v,j}=0\) for each \(j\).
Regard \(\mathcal P_h^k\) as the space of formal forms
\(\sum_i A'_i x_i\), and set \(A=W^\perp\).
Work in the quotient \(\mathcal P_h^k/A_h\), whose evaluation
at \((X_n,x_n)\) is well-defined for large \(n\).

For each \(v\in\mathcal S_f\), take the original basis
\(X^\alpha x_i\) with weights \(\lambda_v\cdot\alpha\).
These are upper bounds for the evaluation slopes since \(x_n\)
is integral. By pairwise coprimality, at \(v\) precisely one generator
\(a_j\) has a negative weight, and the others have weight zero.
Some entry of \(P\) has a term with \(X_j\)-exponent zero,
since \(X_j\) cannot divide all entries of \(P\).
Thus \(K_{\lambda_v,0}(P)=0\) in \eqref{eq:box-filtration}.

At infinity put \(D_i=X^{-\beta_i}V_i\) in every modified row,
where \(X^{\beta_i}\) is the leading monomial of \(V_i\), and put
\(D_i=1\) in every kept row. Thus \(w(D_i)=0\).
Write \(S'=\operatorname{diag}(D_i)S\).
In a modified block, replace \(X^\alpha x_i\) by \(X^\alpha(S'x)_i\)
whenever all terms of the replacement lie in \(\mathcal P_h^k\).
All but \(O(h^{d-1})\) slots of each block admit the replacement,
since the Laurent shifts in \(S'\) are fixed. The resulting forms
constitute a basis. In the modified block the expansion is
triangular with respect to the monomial weight, with a nonzero
pivot at \(X^\alpha x_i\), and the only terms in other blocks
belong to kept coordinates. Give each replaced or kept slot
weight \(w(X^\alpha)+\mu_i\), and each untouched slot in a modified
block weight \(w(X^\alpha)+C'\) for a fixed \(C'>C\).
All these weights bound the evaluations up to \(o(n)\) in the
logarithm. The ambient sum of weights at infinity differs from the
original monomial sum by \(M_h\mathcal B+O(h^{d-1})\).

Denote the weight function from this new basis by \(\mathfrak f_h\)
and the original monomial weight function at infinity by
\(\mathfrak a(A')=\max_i w(A'_i)\).
They differ by \(O(1)\) on nonzero vectors, uniformly in \(h\).
In the forward expansion each term incurs only one of finitely
many shifts. In the inverse expansion solve for a modified
monomial from its replacement and recurse only on lower-weight
monomials in the same block. Along a chain the weights in that
block decrease, and any terms entering kept blocks incur at most
one additional bounded shift. Untouched slots terminate the
recursion. The assigned coordinate costs are also bounded, which
gives the comparison.

On \(A\) consider the additional weight function
\[
 \mathfrak b(A')=\max_i\{\mu_i+w((S'^{-T}A')_i)\}.
\]
It differs from \(\mathfrak a\) by a constant bound, since both the
fixed matrix \(S'^{-T}\) and its inverse have entries of bounded
weight. Choose a nonzero \(D\in\mathcal P\) with \(DS'^{-T}\)
polynomial. For a fixed monomial \(X^\ell\) and a fixed integer \(M\)
large enough, the vectors of \(X^\ell D A_{h-M}\) lie in \(A_h\)
and their transformed coefficients under \(S'^{-T}\) occupy only
replaced or kept slots. Indeed, choose \(\ell\) to pad the lower
faces by more than any fixed replacement shift, and choose \(M\)
to allow for this padding and all upper shifts and polynomial degrees.
For such a vector \(A'\), the equality
\(A'^T x=(S'^{-T}A')^T S'x\) is its expression in the new basis,
so \(\mathfrak f_h=\mathfrak b\) on this subspace.
The codimension is \(O(h^{d-1})\) by Lemma~\ref{lem:box-filtration}.
On the quotient the ordered weights for \(\mathfrak f_h\) and
\(\mathfrak b\) differ by \(O(1)\), so the two weight sums on
\(A_h\) differ by \(O(h^{d-1})\).

To compute the difference of the sums for \(\mathfrak b\) and
\(\mathfrak a\), set \(T=DS'^{-T}\) and use the graph
\[
 G=\{(A',TA'):A'\in A\}\subset F^{2k}.
\]
The joint box \(G_h=G\cap\mathcal P_h^{2k}\) projects onto a space
between \(A_{h-M'}\) and \(A_h\) for a fixed \(M'\), since \(T\)
is polynomial. Its projection thus has codimension \(O(h^{d-1})\)
in \(A_h\). By the uniform weight comparison on the quotient,
the difference of the two weight sums changes by \(O(h^{d-1})\)
upon restriction to this projection.
On the graph, realize \(\mathfrak a\) by costs zero on the first
block and a sufficiently negative constant on the second.
Realize \(\mathfrak b\) by costs \(\mu_i-w(D)\) on the second
block and a sufficiently negative constant on the first.
The fixed bounds for \(T\) and \(T^{-1}\) allow the constants
to be chosen so that the suppressed block has strictly lower
weight on every nonzero element of the graph in each case.

The primitive polynomial Pl\"ucker vector of \(G\) is
\((\bigwedge^c\binom{\mathrm{Id}}{T})P\), whose first-block minors
include the primitive tuple \(P\).
For the costs realizing \(\mathfrak a\), all leading coordinates
of graph vectors are in the first block. The proof of
Lemma~\ref{lem:box-filtration} shows that the leading Pl\"ucker
coordinate uses only these coordinates. Hence the leading
Pl\"ucker weight for these costs is \(\max_I w(P_I)\).
For the costs realizing \(\mathfrak b\), the same argument puts
the leading Pl\"ucker coordinate in the second block, where the
minors are the entries of \(D^c(\bigwedge^c S'^{-T})P\).
The costs cancel the factor \(D^c\) in weight. Also
\(S'^{-T}=\operatorname{diag}(D_i^{-1})S^{-T}\), so each minor
has the same weight as its counterpart for \(S^{-T}\).
The leading Pl\"ucker weight is therefore \(\mathcal K\).
The assertions also hold when \(c=0\), with both weights zero.
Subtracting the two box formulas on \(G_h\) cancels the terms
involving the mean and the multidegrees. Consequently the sum
of the relation weights at infinity is its original sum for
\(\mathfrak a\) plus
\[
 M_h\bigl(\mathcal K-\max_I w(P_I)\bigr)+O(h^{d-1}).
\]

At each place, order the chosen ambient basis by weight, refine
ties by any total order, and row-reduce the relations using the
largest nonzero slots as pivots. The images of the nonpivot
slots are a basis of the quotient. In fact pivot slots can be
eliminated using the relation rows, and no relation is supported
solely on nonpivot slots. The sum of the weights of the chosen
forms is the ambient sum minus the relation sum.
The monomial ambient sums cancel across places by the product
formula. In \eqref{eq:box-filtration}, the \(D_h\overline w_h\)
and multidegree terms cancel for the same reason, and the finite
leading Pl\"ucker weights are zero. The resulting sum of the
weights of all the chosen local forms is
\[
 M_h(\mathcal B-\mathcal K)+O(h^{d-1}).
\]

Choose \(h\) sufficiently large and fixed to make this sum
negative. In any basis of the quotient selected from the
original monomial forms, the evaluations give an integer
vector of sup norm at most exponential in \(n\).
It is nonzero: otherwise all original forms would evaluate to
zero, including the coordinate forms \(x_i\).
The chosen local forms are fixed rational linear forms in these
integer coordinates and are independent at each place.
Their product of absolute values is at most \(e^{-\delta n}\)
for some \(\delta>0\) and sufficiently large \(n\).
The rational integer-point form of the \(S\)-adic Subspace
Theorem \cite[Subspace Theorem]{CZ05} now puts the evaluation
vectors into finitely many proper rational subspaces, since
their sup norms are at most \(e^{C_h n}\) with \(C_h>0\).
Zero values of any of the forms can be included by adjoining
the kernel of that form. A nonzero rational relation in the
quotient basis corresponds to a polynomial form outside
\(A_h\), hence outside \(A\). On an infinite subsequence its
specialization vanishes at \(x_n\), giving a fixed proper
\(F\)-subspace of \(W\).
\end{proof}

\clearpage
\section{The simultaneous cutoff}\label{sec:simcut}

\subsection{Evaluated greatest common divisors}

\begin{lemma}[Evaluated gcd]\label{lem:evaluated-gcd}
Let \(F_1,\ldots,F_t\in\mathbb Z[X_1,\ldots,X_d]\) generate a tuple
of polynomial gcd \(1\) over \(\mathbb Q\).
Then its positive evaluated gcd \(g_n\) satisfies
\(\log g_n=o(n)\). In particular, for a polynomial presentation
\((A_i/B)_{i=1}^m\) with \(B\ne0\) and
\(\gcd(B,A_1,\ldots,A_m)=1\) over \(\mathbb Q\), the least common
positive denominator of its value at \(X_n\) satisfies
\begin{equation}\label{eq:evaluated-denominator}
 \log\operatorname{den}\bigl((A_i/B)(X_n)_{i=1}^m\bigr)
   =nw(B)+o(n).
\end{equation}
\end{lemma}

\begin{proof}
A nonzero constant entry bounds \(g_n\) by its absolute value.
For \(d=1\), a polynomial B\'ezout identity over \(\mathbb Q\),
with denominators cleared, likewise bounds \(g_n\) by a fixed
nonzero integer. Suppose \(d\ge2\) and all nonzero entries are
nonconstant. Fix one nonzero entry \(F_0\). For each irreducible
factor of \(F_0\), the combinations of the tuple divisible by
that factor form a proper linear subspace of the coefficient
space over \(\mathbb Q\). Choose a combination \(G_0\) outside
the finite union of these subspaces, with integral coefficients.
Then \(F_0,G_0\) are coprime over \(\mathbb Q\), and \(g_n\) divides
\(\gcd(F_0(X_n),G_0(X_n))\).

For primes outside \(\mathcal S_f\), the torus gcd theorem of
Corvaja--Zannier and Levin \cite[Theorem~1.13(a)]{LW} bounds
the logarithmic contribution by \(o(n)\).
Indeed, the input coordinates are \(\mathcal S\)-units of heights
\(n\log a_j\), and the exceptional set for each fixed positive
error consists of finitely many translates of proper algebraic
subgroups. Each translate imposes a nontrivial character equation
\((\prod_j a_j^{u_j})^n=b\) with integers \(u_j\) not all zero.
The base is positive and different from \(1\) by pairwise
coprimality, so each translate contains at most one orbit point.

For a prime \(p\in\mathcal S_f\), let \(p\mid a_j\).
Some tuple entry \(H\) is not divisible by \(X_j\), so
\(H_0=H|_{X_j=0}\ne0\). Its evaluated monomials are \(b_1^n,\ldots,b_l^n\)
with distinct positive integers \(b_i\) prime to \(p\).
We show \(\operatorname{ord}_p H_0(X_n)=o(n)\).
For \(l=1\) the valuation is constant. Otherwise use the \(S\)-adic
Subspace Theorem on \((b_1^n,\ldots,b_l^n)\), with places
\(\infty,p\) and all primes dividing the \(b_i\). Take coordinate
forms at every place except \(p\), where one coordinate is replaced
by the linear combination for \(H_0\) with nonzero coefficient on
that coordinate. The forms are independent at each place.
By the product formula and \(|b_i|_p=1\), their product of
absolute values is exactly \(|H_0(X_n)|_p\).
If this were exponentially small infinitely often, the Subspace
Theorem would give a fixed nonzero rational linear relation
among the \(b_i^n\) infinitely often. Its largest base with
nonzero coefficient excludes such a relation.
Now
\[
 H(X_n)\equiv H_0(X_n)\pmod{p^{n\operatorname{ord}_p(a_j)}} ,
\]
so their valuations agree eventually, and the \(p\)-contribution
to \(g_n\) is \(o(n)\). This holds for all the finitely many
generator primes, proving the gcd estimate at all places.
For \eqref{eq:evaluated-denominator}, clear fixed coefficient
denominators in the tuple. For this integral presentation the common denominator is exactly
\[
 \frac{|B(X_n)|}{\gcd(B(X_n),A_1(X_n),\ldots,A_m(X_n))}.
\]
Now apply \eqref{eq:orbit-weight} and the gcd estimate.
\end{proof}


\subsection{The two cutoff criteria}

Fix \(m\ge1\) and \(R=(U_i/V)_{i=1}^m\in F^m\), where
\(U_i,V\in\mathcal P\) and \(V\ne0\). Write \(v=w(V)\) and
\(R_n=R(X_n)\). Denote by \(\operatorname{den}(z)\) the least
common positive denominator of \(z\in\mathbb Q^m\).

\begin{definition}
A number \(Q\ge1\) is \emph{admissible} for \(R\) if, for every
\(\epsilon>0\) and all sufficiently large \(n\),
\begin{equation}\label{eq:admissible}
 q^{1+1/m}\max_{1\le i\le m}|R_{i,n}-p_i/q|\ge e^{-\epsilon n}
 \quad(p\in\mathbb Z^m,\ q\in\mathbb Z,\ 1\le q\le Q^n).
\end{equation}
\end{definition}

For a fixed vector \(Z=(A_i/B)_{i=1}^m\in F^m\), take a polynomial
presentation with \(B\ne0\) and \(\gcd(B,A_1,\ldots,A_m)=1\).
Let \(E_i=U_iB-VA_i\) and \(q_Z(n)=\operatorname{den}(Z(X_n))\).
Define
\begin{equation}\label{eq:line-cutoff}
 s_{\mathrm{line}}=\min_Z
   \left\{w(B):\ \frac{w(B)}{m}+\max_i w(E_i)<v\right\}.
\end{equation}
The set is nonempty since \(Z=R\) qualifies, and the corresponding
value is at most \(v\) after reducing the presentation of \(R\).

To define the subspace criterion number the coordinates of
\((q,p)\) from \(0\) to \(m\) and use the change of coordinates
\[
 S(q,p)=(q,p-qR).
\]
Let \(W\subset F^{m+1}\) have dimension between \(2\) and \(m\), let
\(c=m+1-\dim_F W\), and let \(P\) be the primitive polynomial
Pl\"ucker vector of \(W^\perp\). On coefficient vectors \(S^{-T}\)
adds \(\sum_i R_i A'_i\) to the \(0\)-th coordinate of \(A'\) and
keeps the other coordinates. For \(|I|=c\) with \(0\in I\) the
entries \(C_I=V(\bigwedge^c S^{-T}P)_I\) are polynomial because
only the \(0\)-th coordinate has changed. For \(0\notin I\) the
transformed minor is \(P_I\).
Set
\begin{equation}\label{eq:subspace-weights}
 a_c(W)=\max_{\substack{|I|=c\\0\notin I}} w(P_I),\qquad
 b_c(W)=\max_{\substack{|I|=c\\0\in I}} w(C_I)-v .
\end{equation}
Restrict to subspaces on which \(q\) is not identically zero.
Equivalently, the projection of \(W^\perp\) onto coordinates
\(1,\ldots,m\) is injective, so \(a_c(W)\) is finite.
Define
\begin{equation}\label{eq:subspace-cutoff}
 s_{\mathrm{sub}}=\min_W
 \left\{\frac m c\,a_c(W):\
   \frac m c\,a_c(W)<-\frac{m}{m-c+1}\,b_c(W)\right\}.
\end{equation}
An empty minimum is \(+\infty\), in particular when \(m=1\).
If all \(C_I\) vanish, the right endpoint in the strict
inequality is \(+\infty\).
The minima in \eqref{eq:line-cutoff} and \eqref{eq:subspace-cutoff},
when finite, are attained: polynomial weights are logarithms of
positive integers and \(c\) ranges over a finite set, so there
are only finitely many candidate values in any bounded interval.

\begin{theorem}[Sharp simultaneous cutoff]\label{thm:simcut}
Put \(s_*=\min(s_{\mathrm{line}},s_{\mathrm{sub}})\) and
\(Q_*=\exp(s_*)\).
Every \(1\le Q<Q_*\) is admissible and every \(Q>Q_*\) fails.
If \(s_*<s_{\mathrm{line}}\), the endpoint \(Q_*\) is admissible.
If \(s_*=s_{\mathrm{line}}\), there is a unique vector \(Z_*\)
attaining \eqref{eq:line-cutoff}. In this case \(Q_*\) is admissible
if and only if \(q_{Z_*}(n)>Q_*^n\) for all sufficiently large \(n\).
\end{theorem}

\begin{proof}
For any fixed \(Z\ne R\), Lemma~\ref{lem:evaluated-gcd} and
\eqref{eq:orbit-weight} give
\begin{multline}\label{eq:line-error}
 \log\left(q_Z(n)^{1+1/m}\max_i|R_{i,n}-Z_i(X_n)|\right)\\
   =n\left(\frac{w(B)}m+\max_i w(E_i)-v\right)+o(n).
\end{multline}
For \(Z=R\) the approximation error vanishes. Thus every vector satisfying
\eqref{eq:line-cutoff} gives a counterexample with some fixed
positive \(\epsilon\) at each sufficiently large level where
\(q_Z(n)\le Q^n\).

Suppose now that \eqref{eq:admissible} fails for a fixed
\(\epsilon>0\) at infinitely many levels. Pass to a subsequence
on which \(n^{-1}\log q\to s\in[0,\log Q]\), and choose a fixed
\(F\)-subspace \(W\) of smallest dimension whose specializations
contain infinitely many of the corresponding points \((q,p)\).
Restrict to those points. They have exponentially bounded
coordinates by \eqref{eq:orbit-weight} and the error inequality.
Apply Theorem~\ref{thm:descent} with
\(\mu_0=s\) and \(\mu_i=-s/m-\epsilon\) for \(1\le i\le m\).
Minimality gives \(\mathcal K\le\mathcal B=-m\epsilon\),
so \(W\) is proper because the full space has \(\mathcal K=0\).

If \(\dim_F W=1\), its generator has nonzero \(0\)-th
coordinate since \(q>0\) on the sequence. The ratios \(p_i/q\)
are the specializations \(Z_i(X_n)\) of a fixed \(Z\).
Here \(q_Z(n)\mid q\), and \eqref{eq:evaluated-denominator}
gives \(w(B)\le s\). Multiplying the approximation error
by \(q_Z(n)^{1+1/m}\) still gives a bound by \(e^{-\epsilon n}\).
By \eqref{eq:line-error}, or by exactness when \(Z=R\),
the strict condition in \eqref{eq:line-cutoff} holds.

If \(2\le\dim_F W\le m\), let \(\mathcal K_0(s)\) be the maximum
in \eqref{eq:descent-K} with the same \(s\) and with \(\epsilon=0\).
Every summand changes by at most \(c\epsilon\). Consequently
\begin{equation}\label{eq:subspace-slope}
 \mathcal K_0(s)
   =\max\left\{a_c(W)-\frac{c}{m}s,\
                b_c(W)+\frac{m-c+1}{m}s\right\}
   \le -(m-c)\epsilon <0 .
\end{equation}
In particular \(s\) is strictly inside the interval with
endpoints in \eqref{eq:subspace-cutoff}.
These two cases exclude failure whenever \(\log Q<s_*\).
At \(Q=Q_*\), any infinite sequence of violations must have a
subsequence for which the minimal exceptional subspace is a line.

Conversely, let \(W\) satisfy \eqref{eq:subspace-cutoff} and
take \(s\) strictly inside its interval. Then \(s>0\), since
\(a_c(W)\ge0\). Choose \(\eta>0\) so small that
\(\mathcal K_0(s)+m\eta<0\). For all sufficiently large \(n\)
put \(W_n=W(X_n)\) and rescale by
\[
 T_n(q,p)=\big(e^{-sn}q,\ e^{(s/m+\eta)n}(p-qR_n)\big).
\]
The lattice \(W_n\cap\mathbb Z^{m+1}\) is saturated.
A basis of it extends to a basis of \(\mathbb Z^{m+1}\), since the
quotient is torsion-free, so its wedge is a primitive integral
Pl\"ucker vector. After scaling \(P\) to have integral coefficients,
this wedge is the signed complementary vector of \(P(X_n)\)
divided by the positive gcd of its integer entries, up to an
overall sign. Its Euclidean norm is the covolume in \(W_n\).
Under \(T_n\), the complementary-minor identity shows that
the new covolume equals
\[
 |\det T_n|\,
 \frac{\big\|(\bigwedge^c T_n^{-T})P(X_n)\big\|_2}
      {\gcd_I P_I(X_n)} .
\]
In particular the gcd contributes a factor at most \(1\);
by Lemma~\ref{lem:evaluated-gcd} its logarithm is also \(o(n)\).
The determinant has logarithmic slope \(m\eta\), and inverse
scaling of the error coordinates only decreases the exponents
of the annihilator minors from their values at \(\eta=0\).
By \eqref{eq:orbit-weight} and \eqref{eq:subspace-slope} the
covolume is at most \(\exp(n(\mathcal K_0(s)+m\eta)+O(1))\).
Minkowski's convex body theorem in the lattice span gives a
nonzero point in the unit cube for all sufficiently large \(n\):
the section of the cube contains the Euclidean unit ball in
that span, of volume bounded below in the given dimension.
The point cannot come from \(q=0\), since then each integer
\(p_i\) would have absolute value at most \(e^{-(s/m+\eta)n}<1\).
Changing sign if necessary gives \(1\le q\le e^{sn}\) and
\[
 q^{1+1/m}\max_i|R_{i,n}-p_i/q|
     =q^{1/m}\max_i|qR_{i,n}-p_i|\le e^{-\eta n}.
\]
This proves failure when \(\log Q>s_{\mathrm{sub}}\) by choosing
\(s<\log Q\) inside the interval of a minimizing subspace.
If \(\log Q>s_{\mathrm{line}}\), a minimizing vector gives
failure by \eqref{eq:evaluated-denominator} and
\eqref{eq:line-error}. Together these cover every \(Q>Q_*\).

For uniqueness, suppose two distinct vectors satisfy
\eqref{eq:line-cutoff} at its minimum. Their denominators have
the same leading monomial by distinctness of monomial weights.
Scale the two tuples by nonzero rational constants to cancel this
term upon subtraction. If the difference denominator is nonzero,
it has smaller weight, and each cross-multiplied error of the
difference tuple has weight at most the maximum of the old error
weights. Removing the tuple's polynomial gcd weakly decreases
both the denominator and error weights. The resulting vector
would satisfy \eqref{eq:line-cutoff} at a smaller weight.
If the difference denominator vanishes, a nonzero difference
numerator gives an error divisible by \(V\) of weight at least
\(v\). Both old error weights are strictly less than \(v\),
a contradiction. Thus \(Z_*\) is unique. A counterexample
sequence at \(Q=Q_*\) can only come from such a minimizing
line, by the strict subspace bound and \(w(B)\le s\le s_*\).
Every occurrence of this line has \(q_{Z_*}(n)\mid q\).
Conversely, if \(q_{Z_*}(n)\le Q_*^n\) infinitely often, the
line gives failure as above. This proves both endpoint claims.
\end{proof}

\begin{corollary}[Integral form of the endpoint test]\label{cor:cutoff-gcd}
Suppose \(s_*=s_{\mathrm{line}}\), and choose the tuple for \(Z_*\)
with integer coefficients of joint content \(1\) and positive leading
coefficient \(h\) of \(B\). Let
\(g_n=\gcd(B(X_n),A_1(X_n),\ldots,A_m(X_n))>0\).
For all sufficiently large \(n\), the condition
\(q_{Z_*}(n)\le Q_*^n\) is equivalent to \(g_n\ge h+1\) if the
highest lower term of \(B\) has positive coefficient, and to
\(g_n\ge h\) if it has negative coefficient or there is no lower term.
\end{corollary}

\begin{proof}
The ratio \(B(X_n)/Q_*^n\) is eventually positive and tends to \(h\).
Its difference from \(h\) is eventually of the sign of the highest
lower term, or zero in the absence of lower terms.
The endpoint condition is \(g_n\ge B(X_n)/Q_*^n\).
Since \(h\) and \(g_n\) are integers, the stated alternatives follow.
\end{proof}

\begin{remark}[Bounded supports]\label{rem:cutoff-supports}
The upper bound \(s_{\mathrm{line}}\le v\) restricts the supports
needed for both criteria. In \eqref{eq:line-cutoff} it bounds
\(w(B)\), and the equations \(VA_i=U_iB-E_i\) bound each \(w(A_i)\).
For a subspace candidate at most \(v\), the pure minors have weight
at most \(cv/m\) and the \(C_I\) have weight at most \(v\).
Each mixed minor \(P_I\) has bounded weight because \(VP_I\)
is \(C_I\) minus a linear combination of the pure minors with
coefficients \(\pm U_i\).
When \(m=2\), one has \(c=1\) and these tests reduce to finitely
many rational linear systems on the bounded supports with specified
leading weights. A nonconstant common factor of a feasible tuple
can be removed, lowering the associated threshold while preserving
the strict test. In higher rank the minors must also satisfy the
Pl\"ucker relations over \(F\).
\end{remark}

\begin{remark}[A plane obstruction]\label{rem:plane-obstruction}
Take \(X_n=(2^n,3^n)\), with variables \(X,Y\), and use
\[
 R_1=\frac{Y-1}{X-1},\qquad R_2=\frac{2Y-1}{X-1}.
\]
The plane \(p_2-2p_1=0\) has \(a_1=0\) and \(C_{\{0\}}=1\),
so its interval in \eqref{eq:subspace-cutoff} is
\(0<s<\log2\). A vector with \(w(B)<\log2\) would have constant
denominator, and its first component could not approximate \(R_1\)
with exponentially small error since \(R_1\) has leading base
\(3/2\), different from the positive integer leading bases of
nonzero polynomials in \(X,Y\). Thus \(s_{\mathrm{line}}>0\).
Theorem~\ref{thm:simcut} gives admissibility exactly at \(Q=1\).
\end{remark}

\begin{proposition}[Optimal exponent]\label{prop:optimal-exponent}
If \(Q>1\) is admissible, the estimate \eqref{eq:admissible}
fails with \(1+1/m\) replaced by any \(\kappa<1+1/m\).
\end{proposition}

\begin{proof}
Set \(L=\lfloor Q^{n/m}\rfloor\) and partition the unit \(m\)-cube
into \(L^m\) subcubes of side \(1/L\). Two of the \(L^m+1\)
points \(jR_n\bmod\mathbb Z^m\), \(0\le j\le L^m\), lie in the
same subcube. Their difference gives \(1\le q\le L^m\le Q^n\)
and \(p\in\mathbb Z^m\) with \(\max_i|qR_{i,n}-p_i|\le L^{-1}\).
Thus \(q^{\kappa-1}\max_i|qR_{i,n}-p_i|\) is at most
\(L^{-1+\max(0,m(\kappa-1))}\), which decays exponentially in \(n\).
\end{proof}


\section{Ratios of shifted powers}\label{sec:shifted-powers}

Let \(2\le a<b_1<\cdots<b_m\) be pairwise coprime, with \(m\ge1\).
Put
\begin{equation}\label{eq:shifted-ratios}
 R_n=(R_{i,n})_{i=1}^m,\qquad
 R_{i,n}=\frac{b_i^n-1}{a^n-1},\qquad d_n=\operatorname{den}(R_n).
\end{equation}
The exact common denominator is
\begin{equation}\label{eq:shifted-denominator}
 d_n=\frac{a^n-1}{\gcd(a^n-1,b_1^n-1,\ldots,b_m^n-1)}
     =\exp\bigl(n\log a+o(n)\bigr).
\end{equation}
The asymptotic follows from Lemma~\ref{lem:evaluated-gcd} applied
to \(X-1,Y_1-1,\ldots,Y_m-1\), where \(X,Y_i\) are the variables
for \(a,b_i\). It also follows from the two-base bound of
Bugeaud--Corvaja--Zannier \cite{BCZ}, because the gcd in the
denominator of \eqref{eq:shifted-denominator} divides the gcd
of the first two entries.

\begin{theorem}[Weighted primal estimate]\label{thm:primal}
For every \(\epsilon>0\) and all sufficiently large \(n\), for
every integer \(1\le q<d_n\) and every \(t_i\ge0\) with
\(\sum_i t_i=n^{-1}\log q\), one has
\begin{equation}\label{eq:weighted-primal}
 \max_{1\le i\le m} e^{nt_i}\|qR_{i,n}\|\ge e^{-\epsilon n},
\end{equation}
where \(\|u\|\) denotes distance to the nearest integer.
In particular the cutoff in \eqref{eq:admissible} is \(Q_*=a\),
with failure at equality.
\end{theorem}

\begin{proof}
Suppose that \eqref{eq:weighted-primal} fails infinitely often for a
fixed \(\epsilon>0\). Choose nearest integers \(p_i\) and pass to a
subsequence on which \(n^{-1}\log q\to s\) and the weight parameters
converge to \(t_1,\ldots,t_m\), using the same notation for their
limits. These satisfy \(0\le t_i\), \(\sum_i t_i=s\le v=\log a\).
Here \(V=X-1\), \(U_i=Y_i-1\).
As in the cutoff proof take a fixed minimal-dimensional \(W\)
containing infinitely many of the points \((q,p)\) upon specialization
and restrict to those points. The coordinates are exponentially
bounded since \(q<d_n\) and the \(p_i\) are nearest to \(qR_{i,n}\).
Use \(\mu_0=s\) and \(\mu_i=-t_i-\epsilon\) in
Theorem~\ref{thm:descent}. Minimality gives
\(\mathcal K\le\mathcal B=-m\epsilon\), so \(W\) is proper.

If \(W\) has dimension at least two, then \(1\le c\le m-1\).
For a pure minor \(P_I\), with \(0\notin I\), the inequality
\(\mathcal K\le-m\epsilon\) gives
\[
 w(P_I)\le\sum_{i\in I}t_i-(m-c)\epsilon<v .
\]
Every pure minor is constant because all nonconstant monomials
have weight at least \(v\). At least one such \(P_I\) is nonzero
since \(q\) is nonzero on the sequence. Choose \(j\in I\) and
put \(J=\{0\}\cup(I\setminus\{j\})\).
In the mixed numerator \(C_J\) of \eqref{eq:subspace-weights},
setting \(X=1\) kills the term \(VP_J\); what remains is a
constant linear combination of the \(Y_i-1\), with nonzero
\(Y_j\)-coefficient \(\pm P_I\). Thus \(w(C_J)\ge\log b_j>v\),
since setting \(X=1\) cannot increase the weight.
The contribution of \(J\) to \(\mathcal K\) is
\[
 w(C_J)-v+s-\sum_{i\in J\setminus\{0\}}t_i-(c-1)\epsilon
   >-(c-1)\epsilon,
\]
contradicting \(\mathcal K\le-m\epsilon\).

If \(W\) has dimension one, \(p_i/q=Z_i(X_n)\) for a fixed
\(Z=(A_i/B)_{i=1}^m\) in the reduced polynomial presentation
of \eqref{eq:line-cutoff}. Since \(q_Z(n)\mid q\),
Lemma~\ref{lem:evaluated-gcd} gives \(w(B)\le s\le v\).
Thus \(B\) is constant or linear in \(X\), with rational
coefficients. For each nonzero \(E_i=U_iB-VA_i\),
\eqref{eq:orbit-weight} gives
\[
 \log|q(R_{i,n}-Z_i(X_n))|
       =n(s+w(E_i)-v-w(B))+o(n).
\]
Its slope is negative by the assumed bound with \(t_i\ge0\).
Since \(s\ge w(B)\), this forces \(w(E_i)<v\). Hence each
\(E_i\) is constant or zero. The identity
\(E_i|_{X=1}=B(1)(Y_i-1)\) gives \(B(1)=0\), then every
\(E_i=0\). Thus \(Z=R\), which implies \(d_n\mid q\) and
contradicts \(q<d_n\).

Taking all \(t_i=(\log q)/(mn)\) gives the estimate for
\eqref{eq:admissible} throughout \(1\le q<d_n\).
By \eqref{eq:shifted-denominator}, \(Q^n<d_n\) eventually
for every \(Q<a\). Exact approximation at \(q=d_n\le a^n-1\)
proves failure at \(Q=a\) and above.
\end{proof}

\begin{theorem}[Dual estimate and cutoff]\label{thm:dual}
Write \(H=\max_i|k_i|\) for nonzero \(k\in\mathbb Z^m\).
For every \(\epsilon>0\) and all sufficiently large \(n\), each
such \(k\) with \(k\cdot R_n\notin\mathbb Z\) satisfies
\begin{equation}\label{eq:dual-estimate}
 H^m\|k\cdot R_n\|\ge e^{-\epsilon n}.
\end{equation}
The smallest \(H\) of a nonzero \(k\) with \(k\cdot R_n\in\mathbb Z\)
is \(\exp(n(\log a)/m+o(n))\). With exact values included,
\eqref{eq:dual-estimate} holds for all \(1\le H\le Q^n\) at
each \(1\le Q<a^{1/m}\), and fails at \(Q=a^{1/m}\) and above.
\end{theorem}

\begin{proof}
The lattice
\(\Lambda_n=\{(q,p-qR_n): (q,p)\in\mathbb Z^{m+1}\}\)
has covolume one. For \(1\le T<d_n\) scale it by
\(\operatorname{diag}(T^{-1},T^{1/m},\ldots,T^{1/m})\).
Fix \(\eta>0\). Uniformly in \(T\) and all sufficiently large
\(n\), the first successive minimum of the scaled lattice
in the unit cube is at least \(e^{-\eta n}\).
Indeed, a vector with \(|q|>T\) has first coordinate of
absolute value greater than \(1\).
If \(1\le |q|\le T\), use \eqref{eq:weighted-primal} with
equal weights for \(|q|\) and use \(T^{1/m}\ge|q|^{1/m}\).
At \(q=0\), a nonzero vector has an integer residual of
sup norm at least \(1\) before scaling.
The scaled lattice still has covolume one. Minkowski's second
theorem therefore bounds the last minimum by
\(C_m e^{m\eta n}\) for a positive constant \(C_m\).
A nonzero dual lattice vector has a nonzero integral pairing
with at least one of \(m+1\) independent primal lattice vectors
of sup norm at most \(C_m e^{m\eta n}\). Its sup norm is
therefore at least \(c_m e^{-m\eta n}\) for some \(c_m>0\).
The dual vectors are precisely
\[
 \bigl(T(\ell+k\cdot R_n),T^{-1/m}k\bigr),
       \qquad (\ell,k)\in\mathbb Z^{m+1}.
\]

Set \(v=\log a\) and choose \(0<\delta<\min(v/m,\epsilon/m)\).
For \(H\le e^{(v/m-\delta)n}\), take
\(T=H^m e^{(m^2+1)\eta n}\) with \(\eta>0\) so small that
\((m^2+1)\eta<m\delta\) and
\((m^2+m+1)\eta<\epsilon\).
By \eqref{eq:shifted-denominator}, \(1\le T<d_n\) eventually.
The last \(m\) coordinates of the dual vector have sup norm
\(e^{-(m+1/m)\eta n}<c_m e^{-m\eta n}\) for all large \(n\).
Choosing \(\ell\) to minimize its first coordinate now gives
\[
 H^m\|k\cdot R_n\|\ge c_m e^{-(m^2+m+1)\eta n}.
\]
This both proves \eqref{eq:dual-estimate} in this range and
excludes exact values. For \(H>e^{(v/m-\delta)n}\), a nonexact
value is at least \(1/d_n\) from the nearest integer. Since
\(d_n\le e^{vn}\), this yields \(H^m\|k\cdot R_n\|>e^{-m\delta n}\).

The lattice of \(k\in\mathbb Z^m\) with \(k\cdot R_n\in\mathbb Z\)
has index \(d_n\), since the image of \(k\mapsto k\cdot R_n\bmod
\mathbb Z\) is the cyclic subgroup of order \(d_n\)
generated by the coordinates of \(R_n\). The cube of sup norm
at most \(a^{n/m}\) has volume \(2^m a^n>2^m d_n\).
Minkowski's convex body theorem gives a nonzero vector of this
lattice in the cube. Together with the exclusion of exact
values for every fixed \(\delta>0\) in the smaller range, this
proves the asymptotic for the least \(H\) and failure at
\(Q=a^{1/m}\). For a fixed \(Q<a^{1/m}\) one can take \(\delta\)
small enough that \(\log Q<v/m-\delta\).
\end{proof}

\begin{corollary}\label{cor:binary-ternary-quinary}
For \(R_n=((3^n-1)/(2^n-1),(5^n-1)/(2^n-1))\), the sharp
primal exponent is \(3/2\), with cutoff \(2\), and the dual
coefficient cutoff is \(\sqrt{2}\). Each cutoff fails at equality.
\end{corollary}

\begin{proof}
Apply Theorems~\ref{thm:primal} and~\ref{thm:dual} with
\((a,b_1,b_2)=(2,3,5)\), and Proposition~\ref{prop:optimal-exponent}
with any \(1<Q<2\).
\end{proof}

\begin{proposition}[Weighted covering]\label{prop:weighted-covering}
Fix \(\eta>0\). For all sufficiently large \(n\), let
\(t_i\ge0\), \(\sum_i t_i=s\), and \(1\le e^{sn}<d_n\).
For every \(q_0\in\mathbb R\) and \(y\in\mathbb R^m\), there is
\(q\in\mathbb Z\) such that
\[
 |q-q_0|\ll_m e^{(s+m\eta)n},\qquad
 \|qR_{i,n}-y_i\|\ll_m e^{(-t_i+m\eta)n}\quad(1\le i\le m).
\]
The constants are uniform in the weights and shifts.
\end{proposition}

\begin{proof}
Scale \(\Lambda_n\) by
\(\operatorname{diag}(e^{-sn},e^{nt_1},\ldots,e^{nt_m})\),
a map of determinant one.
For \(1\le |q|\le e^{sn}\) and \(s>0\), apply the weighted
primal bound with \(t_i\) multiplied by \(\log|q|/(ns)\).
These new weights are no larger and sum to \(n^{-1}\log|q|\).
When \(s=0\) and \(|q|=1\), the weights are already zero.
Vectors with \(|q|>e^{sn}\) or with \(q=0\) and \(p\ne0\)
have sup norm at least one after scaling.
Thus the first minimum in the cube is at least \(e^{-\eta n}\).
As in Theorem~\ref{thm:dual}, the last minimum is
\(O_m(e^{m\eta n})\). A parallelepiped spanned by \(m+1\)
independent lattice vectors of this size covers the space
under translations by the sublattice they generate, so the
scaled lattice has covering radius \(O_m(e^{m\eta n})\)
in the sup norm. Apply this bound to the image of
\((q_0,-y)\), then undo the scaling.
\end{proof}

\section{The scalar exponent and its sharp range}\label{sec:scalar}
Use all primes occurring in the roots as orbit generators \(a_j\), so every term corresponds to a monomial in \(X_j=a_j^n\). For constants take one auxiliary prime if necessary. Write \(R=U/V\) in reduced form over \(\mathbb Q[X_1,\ldots,X_d]\). For a polynomial \(P\ne0\) let \(L(P)=\exp w(P)\), its largest integer monomial base, and set \(L(0)=0\).
\begin{theorem}\label{thm:scalarcut}
The sharp scalar cutoff, with exponent two, is the positive integer
\begin{equation}\label{eq:scalarmin}
   Q_*=\min_{\substack{Z=A/B\\ B\ne0,\ (A,B)=1}}
      \{L(B):\ L(B)L(UB-VA)<L(V)\}.
\end{equation}
It is at most \(L(V)\) and the minimizing function \(Z_*\) is unique. The endpoint \(Q_*\) is admissible if and only if \(\operatorname{den}Z_*(X_n)>Q_*^n\) for all sufficiently large \(n\). A finite list of rational linear feasibility problems determines \(Q_*\) and \(Z_*\).
\end{theorem}
\begin{proof}
In dimension one the exact function is a candidate and the only proper subspaces of the two-dimensional coefficient space that can contain an approximant are lines. Thus Theorem~\ref{thm:simcut} gives \eqref{eq:scalarmin} and the stated endpoint test. More explicitly, if \(UB-VA\ne0\), the exponential growth rate of
\(\operatorname{den}(Z(X_n))^2|R(X_n)-Z(X_n)|\)
is the logarithm of \(L(B)L(UB-VA)/L(V)\). Distinct bad fractions of the same denominator growth rate cannot both approximate \(R\) this well: their distance at every large level is at least the reciprocal of the product of their reduced denominators.

Enumerate the finitely many monomial bases \(b\le L(V)\) in increasing order. In \(B\) allow monomials of base at most \(b\) and set its coefficient of base \(b\) to one. In \(A\) allow bases at most \(b\max(1,L(U)/L(V))\). This suffices since a qualifying function eventually has \(|A(X_n)/B(X_n)|\le |R(X_n)|+1\). Set to zero the coefficients in \(UB-VA\) of all monomials of base \(m\) with \(bm\ge L(V)\). These are affine linear equations over \(\mathbb Q\). A common nonconstant factor of a feasible pair can be cancelled, decreasing \(b\) and preserving feasibility at the reduced denominator base. The first consistent system therefore supplies the reduced minimizer.
\end{proof}

\begin{proof}[Proof of Theorem~\ref{thm:cz}]
A denominator of base one is constant. A candidate of base one would give a polynomial \(\zeta\) in the prime powers with \(R(X_n)-\zeta(X_n)\) tending to zero. This contradicts \eqref{eq:czhyp}, which is unchanged by multiplying or dividing both representatives by their common polynomial factor. Thus \(Q_*\ge2\). The function \((3/2)^n\) satisfies \eqref{eq:czhyp}, and its exact denominator \(2^n\) gives failure at \(Q=2\).
For optimality of the exponent, Dirichlet's theorem with \(K=\lfloor Q^n\rfloor\) gives \(1\le q\le K\) with \(q|R_n-p/q|\le1/K\). For each \(k<2\) this implies
\(q^k|R_n-p/q|\le K^{-\min(1,2-k)}\),
which violates the corresponding lower bound for a sufficiently small fixed \(\epsilon>0\).
\end{proof}

For \(7^n/(2^n+1)\) and \(7^n/(2^n-1)\) the cutoff equals two: the leading base \(7/2\) excludes polynomial approximation and the exact denominator polynomial has base two. In the first case \(7\nmid 2^n+1\), since \(2\) has odd order modulo \(7\), so the endpoint holds. In the second case the reduced denominator is at most \(2^n-1\), giving failure. An earlier \emph{nonexact} approximant can also decide the cutoff. For \(R_n=(15^n+2^n)/10^n\) the reduced denominator is \(10^n\), whereas \(Z_n=3^n/2^n\) gives \(2^{2n}|R_n-Z_n|=(4/5)^n\). Here again the leading base \(3/2\) excludes denominator base one.

\begin{proposition}\label{prop:exceptions}
With the roots, the number of terms, \(1<Q<2\), and \(\epsilon>0\) fixed, the number of exceptional indices in Theorem~\ref{thm:cz} is unbounded as the coefficients vary. The exceptional rational approximants can be distinct and nonexact.
\end{proposition}
\begin{proof}
For any \(L\ge1\), take an integer \(C>((3/2)e^\epsilon)^L\) and put \(u(n)=C4^n+3^n,\ v(n)=C2^n\). In \(u(n)/v(n)=2^n+C^{-1}(3/2)^n\), no integer-root power sum cancels the term of base \(3/2\), so \eqref{eq:czhyp} holds. For every \(1\le n\le L\), the integer \(p=2^n\) has positive error less than \(e^{-\epsilon n}\). In the lifted coordinates these violations already satisfy one fixed relation, \(p=2^nq\). A bound on the number of exceptional subspaces, as in the quantitative Subspace Theorem \cite{EF}, therefore does not by itself bound the number of indices.
\end{proof}

\section{Scalar approximation at a fixed embedding}
\label{sec:algscalar}

Let \(u(n)\) and \(v(n)\) be finite sums of terms \(P(n)\gamma^n\),
where \(P\) has algebraic coefficients and \(\gamma\) is a nonzero
algebraic number. Such terms include the contributions of repeated
characteristic roots of a linear recurrence. A zero characteristic
root contributes only to an initial segment, which can be discarded.
Fix an embedding of the data in \(\mathbb C\), and set \(R_n=u(n)/v(n)\)
at the levels where \(v(n)\ne0\). All errors below use the ordinary
complex modulus at this embedding.

\begin{definition}
A real number \(Q\ge1\) is admissible if, for every \(\epsilon>0\),
all sufficiently large defined levels satisfy
\begin{equation}\label{eq:alg-admissibility}
 q^2|R_n-p/q|\ge \exp(-\epsilon n)
 \qquad(p,q\in\mathbb Z,\ 1\le q\le Q^n).
\end{equation}
For \(z\in\mathbb Q\), let \(\operatorname{den}(z)\) be its positive
reduced denominator, with \(\operatorname{den}(0)=1\).
\end{definition}

It suffices to test reduced fractions, since reducing preserves the
target \(p/q\) and decreases the left side. For real \(R_n\) with
infinitely many defined levels, the exponent \(2\) is optimal whenever
\(Q>1\) is admissible. Indeed, with \(J=\lfloor Q^n\rfloor\),
Dirichlet's theorem gives \(1\le q\le J\) and
\(q|R_n-p/q|\le J^{-1}\). For any \(\kappa<2\) this yields
\(q^\kappa|R_n-p/q|\le J^{-\min(1,2-\kappa)}\).

\subsection{The character ring and the cutoff}

Identify the prescribed embedding with the inclusion of the data in
\(\mathbb C\). Let \(K\subset\mathbb C\) be a finite Galois extension of
\(\mathbb Q\) containing the data and write
\(G=\operatorname{Gal}(K/\mathbb Q)\). Choose \(m\ge1\) killing the
torsion of the multiplicative group generated by the roots and all
their conjugates. Their \(m\)-th powers generate a \(G\)-stable
torsion-free group \(\Lambda\). For a residue \(0\le r<m\) put
\(n=r+mk\). There are \(U,V\) in the ring
\[
 \mathcal A=K[T][\Lambda],\qquad
 P[k]=\sum_{\gamma}P_\gamma(k)\gamma^k
 \quad\text{for } P=\sum_\gamma P_\gamma(T)[\gamma],
\]
such that \(U[k]=u(r+mk)\) and \(V[k]=v(r+mk)\).
Indeed, the summand \(P(n)\gamma^n\) becomes
\(\gamma^r P(r+mT)[\gamma^m]\). The character notation satisfies
\([\gamma][\delta]=[\gamma\delta]\). The action of \(G\) fixes \(T\) and
sends \(c[\gamma]\) to \(\sigma(c)[\sigma\gamma]\). Evaluation commutes
with this action. Omit residues with \(V=0\).

Let \(M_K\) be the set of places of \(K\). Use absolute values
normalized for the product formula, with local degree divided by
\([K:\mathbb Q]\). At an archimedean place \(v\) represented by
\(\sigma_v:K\hookrightarrow\mathbb C\), write
\(|x|_v=|\sigma_v(x)|^{t_v}\), where \(t_v=1/[K:\mathbb Q]\) for real
places and \(2/[K:\mathbb Q]\) for complex places. Thus
\(\sum_{v\mid\infty}t_v=1\). Fix a finite set \(S\) containing the
archimedean places and every finite place at which some basis element
of \(\Lambda\) is a nonunit, and write \(S_f\) for its finite places.
For \(P\in\mathcal A\) let its support be the characters with
\(P_\gamma(T)\ne0\). Define
\[
 w_v(P)=\max_{\gamma\in\operatorname{supp}P}\log|\gamma|_v,\qquad
 \rho(P)=\max_{\gamma\in\operatorname{supp}P}\log|\gamma|.
\]
The weights of zero are \(-\infty\).

For \(Z\in\operatorname{Frac}(\mathcal A)^G\), choose a coprime
representation \(Z=A/B\) in \(\mathcal A\), with \(0=0/1\). In
particular the reduction cancels common factors in \(K[T]\).
For the current residue define
\begin{align}
 D(Z)&=\sum_{v\mid\infty}w_v(B)
       +\sum_{v\in S_f}\max\{w_v(A),w_v(B)\}, \label{eq:alg-D}\\
 e(Z)&=\rho(UB-VA)-\rho(V)-\rho(B). \label{eq:alg-e}
\end{align}
Since \(\mathcal A\) is a Laurent polynomial ring over \(K[T]\), its
units are \(c[\gamma]\) with \(c\in K^\times\). Multiplying both
members of the pair by such a unit leaves \(D\) unchanged by the
product formula, and leaves \(e\) unchanged by cancellation of the
common weight. Enlarging \(S\) leaves \(D\) unchanged. Evaluation of
\(Z\) gives rational numbers wherever defined.

\begin{theorem}\label{thm:algcut}
On each residue with \(V\ne0\), let
\begin{equation}\label{eq:alg-variational}
 d_r=\inf\{D(Z):Z\in\operatorname{Frac}(\mathcal A)^G,\
                       2D(Z)+e(Z)<0\},
 \qquad
 Q_*=\min_{r:\,V\ne0}\exp(d_r/m).
\end{equation}
Use \(+\infty\) for an empty infimum or minimum. Each finite \(d_r\)
is attained by a unique function \(Z_r\).
Every \(1\le Q<Q_*\) is admissible, and every \(Q>Q_*\) is
inadmissible. If \(Q_*<\infty\), then \(Q_*\) is admissible if and
only if, on each residue with \(\exp(d_r/m)=Q_*\),
\begin{equation}\label{eq:alg-endpoint}
 \operatorname{den}(Z_r[k])>Q_*^{\,r+mk}
 \qquad\text{for all sufficiently large }k.
\end{equation}
If every \(d_r\) is infinite, every finite \(Q\ge1\) is
admissible. If there are no defined residues, admissibility is
vacuous.
\end{theorem}

\subsection{Orbit evaluation}

\begin{lemma}\label{lem:alg-orbit}
A nonzero \(P\in\mathcal A\) has only finitely many zeros on the
orbit. If \(Z\in\operatorname{Frac}(\mathcal A)\) has rational values
at infinitely many levels, then \(Z\) is \(G\)-invariant.
\end{lemma}

\begin{proof}
The sequence \(P[k]\) satisfies a linear recurrence. By the
Skolem--Mahler--Lech theorem \cite{LechSML}, infinitely many zeros
would give a full infinite arithmetic progression of zeros. On
\(k=a+b\ell\), \(b\ge1\), the distinct characters in \(P\) have
distinct \(b\)-th powers by torsion-freeness. The corresponding
polynomial multiples of the exponentials in \(\ell\) are linearly
independent. Indeed, if \(E\) is the forward shift, \(E-\lambda\)
lowers polynomial degree on terms with base \(\lambda\) and
preserves degree with nonzero leading coefficient on terms with a
different base. A product of powers of these operators can isolate
any one base of the alleged relation. This contradicts the existence
of the progression. For the second assertion, the cross numerator
of \(Z-\sigma Z\) vanishes at infinitely many levels for every
\(\sigma\in G\) and therefore vanishes in \(\mathcal A\).
\end{proof}

We use the number-field Subspace Theorem \cite{Schlickewei} in
the normalized form of \cite[proof of Lemma~1]{CZ04}. Write
\(\|\mathbf x\|_v=\max_i|x_i|_v\) and
\(H(\mathbf x)=\prod_{v\in M_K}\|\mathbf x\|_v\) for the
multiplicative projective height. For \(N\ge2\), take \(N\) independent
\(K\)-linear forms \(L_{v,i}\) at each \(v\in S\). Nonzero vectors
satisfying
\[
 \prod_{v\in S}\prod_{i=1}^{N}
     \frac{|L_{v,i}(\mathbf x)|_v}{\|\mathbf x\|_v}
     < H(\mathbf x)^{-N-\delta}
 \qquad(\delta>0)
\]
are contained in finitely many proper \(K\)-linear subspaces.
For vectors integral outside \(S\),
\(\prod_{v\in S}\|\mathbf x\|_v\ge H(\mathbf x)\). In particular,
if \(H(\mathbf x)\le \exp(Ck)\) for \(C>0\), the unnormalized form
product being at most \(\exp(-ck)\), with fixed \(c>0\), suffices
by taking \(0<\delta<c/C\).

\begin{lemma}[Local weights]\label{lem:localweight}
For \(P\in\mathcal A\setminus\{0\}\) and any fixed place \(v\),
\begin{equation}\label{eq:alg-localweight}
 \log |P[k]|_v=k w_v(P)+o(k).
\end{equation}
At the specified embedding,
\(\log|P[k]|=k\rho(P)+o(k)\). In particular a nonzero exponential
polynomial whose characters all have ordinary modulus \(1\) at that
embedding has logarithmic modulus \(o(k)\).
\end{lemma}

\begin{proof}
Enlarge \(S\) to include \(v\). Expand \(P\) into its distinct
monomials \(T^j[\gamma]\) and use their values
\(x_{j,\gamma}=k^j\gamma^k\) as coordinates, for \(k\ge1\).
At \(v\) replace the coordinate form of one maximal-weight monomial
by the form evaluating to \(P[k]\). The coefficient of the pivot
monomial is nonzero, so the forms are independent. Use coordinate
forms elsewhere. The product of the original coordinate forms over
\(S\) is at most a fixed power of \(k\) by the product formula for the
characters and \(\prod_{w\in S}|k|_w\le k\). The replacement changes
the product by
\[
 \frac{|P[k]|_v}{|k|_v^j\exp(k w_v(P))}.
\]
At a fixed place \(|k|_v^{-j}\) is bounded by a fixed power of \(k\).
An exponential deficit in \eqref{eq:alg-localweight} along infinitely
many \(k\) would therefore give an exponentially small form product.
The coordinates are integral outside \(S\) and their logarithmic
height is \(O(k)\). The Subspace Theorem would give a fixed nonzero
linear relation among the coordinates at infinitely many \(k\),
contrary to Lemma~\ref{lem:alg-orbit}. With one coordinate the
conclusion follows directly. The upper estimate follows by bounding
the finitely many terms by \(\exp(k w_v(P))k^{O(1)}\), and division by
the local normalization exponent gives the assertion for ordinary
modulus.
\end{proof}

\begin{lemma}\label{lem:alg-slopes}
If \(Z\in\operatorname{Frac}(\mathcal A)^G\), then \(D(Z)\ge0\) and
\[
 \log\operatorname{den}(Z[k])=kD(Z)+o(k).
\]
If \(Z\ne U/V\), then
\(\log|R_{r+mk}-Z[k]|=ke(Z)+o(k)\). For \(Z=U/V\) the error is zero
and \(e(Z)=-\infty\).
\end{lemma}

\begin{proof}
Choose a basis \(\gamma_1,\ldots,\gamma_d\) of \(\Lambda\), so
that \(\mathcal A=K[T,X_1^{\pm1},\ldots,X_d^{\pm1}]\).
Multiply the reduced pair by a common monomial in the \(X_i\) to
make its entries ordinary polynomials with no common coordinate
factor. They are then coprime in \(K[T,X_1,\ldots,X_d]\).
Enlarge \(S\) so their coefficients are \(S\)-integers. For this
pair,
\begin{equation}\label{eq:alg-offS}
 -\sum_{v\notin S}\log\max(|A[k]|_v,|B[k]|_v)=o(k).
\end{equation}
In rank zero this follows from a polynomial B\'ezout identity
over \(K\), after including the coefficient places of the identity.
It is immediate when one component is a nonzero constant that is
an \(S\)-unit. Otherwise, in positive rank apply Xiao's
almost-unit gcd estimate \cite[preprint, Corollary~3.4 and the proof of
Theorem~4.2]{Xiao} to
\(\mathbf y_k=(k,\gamma_1^k,\ldots,\gamma_d^k)\). In Xiao's almost-unit condition the sum of the logarithmic
local heights of \(\mathbf y_k\) and its coordinatewise inverse
outside \(S\) is \(O(\log k)\), since only the coordinate \(k\)
can contribute. The affine height \(h([1:\mathbf y_k])\) and the
largest coordinate height are both comparable to \(k\), since
\(h(\gamma_i^k)=k h(\gamma_i)\) and \(h(\gamma_i)>0\) by
Northcott's theorem and torsion-freeness. Thus the almost-unit
condition holds eventually for every fixed positive tolerance. A proper exceptional closed subset of the torus is cut
out, in particular, by some nonzero Laurent polynomial. Clearing
negative powers of \(T\) and applying Lemma~\ref{lem:alg-orbit}
shows that the orbit eventually avoids it. For every \(\eta>0\)
Xiao's estimate bounds the nonnegative left side of
\eqref{eq:alg-offS} by \(\eta\max_i h(y_{k,i})\) eventually.
Here \(h\) denotes absolute logarithmic height.

Combining \eqref{eq:alg-offS} with Lemma~\ref{lem:localweight}
gives
\[
 h([A[k]:B[k]])=
 k\sum_{v\in S}\max\{w_v(A),w_v(B)\}+o(k).
\]
For \(z\in\mathbb Q\),
\(h([z:1])=\log\operatorname{den}(z)+\log^+|z|\).
Since \(Z[k]\) is rational, \(\log^+|Z[k]|\) equals
\(\sum_{v\mid\infty}\log^+|Z[k]|_v\). Subtracting this sum leaves
the archimedean weights \(w_v(B)\), giving \eqref{eq:alg-D}.
The case \(Z=0\) has denominator one. The assertion about \(e(Z)\)
follows by evaluating \((UB-VA)/(VB)\) at the specified embedding
and applying Lemma~\ref{lem:localweight}.
\end{proof}

\subsection{A finite list at each ceiling and error}

\begin{lemma}\label{lem:alg-finite}
Fix a residue with \(V\ne0\), \(Q\ge1\), and \(\epsilon>0\).
There is a finite set of functions in
\(\operatorname{Frac}(\mathcal A)^G\) such that every violation of
\eqref{eq:alg-admissibility} at a sufficiently large level on the
residue satisfies \(p/q=Z[k]\) for one of those functions.
\end{lemma}

\begin{proof}
Include \(Z=0\) to account for \(p=0\). For \(p\ne0\) consider two
copies, multiplied by \(p\) and by \(q\), of the monomials
\(T^\ell[\gamma^{\mathbf j}]\) evaluated at \(k\). Here
\(\gamma^{\mathbf j}=\prod_{i=1}^d\gamma_i^{j_i}\) and
\((\ell,\mathbf j)\in\{0,\ldots,h\}^{d+1}\), a box with
\(M=(h+1)^{d+1}\) indices.

At each archimedean place set
\(U_v=\sigma_v^{-1}U\) and \(V_v=\sigma_v^{-1}V\), using the
semilinear action on both coefficients and characters.
Since \(K\) is Galois, \(\sigma_v\) is an automorphism of \(K\).
As \(\sigma_v\) commutes with evaluation,
\(\sigma_v(U_v[k]/V_v[k])=R_{r+mk}\) at the prescribed embedding.
Thus the forms at all archimedean places measure the same complex
error.

Order the monomials by local character weight, then by \(T\)-degree,
then lexicographically by the character exponents. This order on
the exponent vectors is translation invariant. Let \(a_v\) be the
largest monomial of \(V_v\) in this order. For each box monomial \(b\)
for which both translated supports of \((b/a_v)U_v\) and
\((b/a_v)V_v\) fit in the box, replace the \(p\)-coordinate form at
\(b\) by the form evaluating to
\[
 (b/a_v)[k]\bigl(p V_v[k]-q U_v[k]\bigr).
\]
Here the monomial \(b/a_v\) may contain negative exponents before
multiplication, and its evaluation is defined for \(k\ge1\).
Keep the other forms and all \(q\)-coordinate forms. The
\(p\)-block is triangular in the given order with nonzero diagonal,
so the forms are independent. The translated supports impose only
fixed-width exclusions at the boundary of the box. There are
\(I_v=M-O(h^d)\) replacements. At every finite place of \(S\) keep
all coordinate forms.

The quotient \(|V_v[k]/a_v[k]|_v\) is \(O(k^C)\) for a fixed \(C\)
because \(a_v\) has maximal local character weight. Thus each
replacement changes the local form product by a factor at most
\(O(k^C)(|p-qR_{r+mk}|/|p|)^{t_v}\).
In the unmodified product over \(S\) the characters cancel by
the product formula, powers of \(k\) cost \(k^{O_h(1)}\), and
\(\prod_{v\in S}|pq|_v\le |p|q\).
Writing \(I=\sum_{v\mid\infty}t_v I_v=M-O(h^d)\), the modified
product is at most
\begin{equation}\label{eq:alg-boxprod}
 O\bigl(k^{C_h}\bigr)(|p|q)^{M-I}
       \bigl(q|p-qR_{r+mk}|\bigr)^I.
\end{equation}
For exceptional fractions, \(\log(|p|q)=O(k)\) with a constant
independent of \(h\), because \(q\le Q^{r+mk}\),
\(|p|\le q|R_{r+mk}|+1\), and the local weight lemma gives
\(|R_{r+mk}|\le\exp(O(k))\). Also
\(q|p-qR_{r+mk}|<\exp(-\epsilon(r+mk))\). Choose \(h\) large
enough that the negative exponential contribution in
\eqref{eq:alg-boxprod} dominates the boundary contribution.
For that fixed \(h\), \eqref{eq:alg-boxprod} is exponentially
small in \(k\), including when the error vanishes.

The coordinate vector is integral outside \(S\), with logarithmic
height \(O_h(k)\). Apply the Subspace Theorem and choose a fixed
nonzero relation for each of the resulting subspaces. Each relation
is \(pB[k]=qA[k]\) for fixed \(A,B\in\mathcal A\) with
\((A,B)\ne(0,0)\). If \(B=0\), Lemma~\ref{lem:alg-orbit} makes
this relation impossible at large levels. Otherwise \(B[k]\ne0\)
eventually. Discard the finitely many levels accounted for by any
relation used only finitely often. Every remaining quotient \(A/B\)
takes rational values at infinitely many levels and hence is
invariant by Lemma~\ref{lem:alg-orbit}.
\end{proof}

\begin{proof}[Proof of Theorem~\ref{thm:algcut}]
Call an invariant function an obstruction on the residue when
\(2D(Z)+e(Z)<0\). For a nonexact function,
Lemma~\ref{lem:alg-slopes} gives
\[
 \log\bigl(\operatorname{den}(Z[k])^2|R_{r+mk}-Z[k]|\bigr)
   = k(2D(Z)+e(Z))+o(k).
\]
Thus a function on the finite list of Lemma~\ref{lem:alg-finite}
can occur in infinitely many violations for fixed \(\epsilon\) only
if it is an obstruction, since \(q\ge\operatorname{den}(Z[k])\).
Conversely each obstruction gives violations at its reduced
denominators for some fixed \(\epsilon>0\) whenever they fit the
ceiling and \(k\) is sufficiently large. Both statements include the
exact function when it is invariant.

To show attainment, fix an upper bound \(D_0\) for the slopes of
interest. Choose \(C\ge0\) such that
\(\log^+|R_{r+mk}|\le Ck+o(k)\), using
Lemma~\ref{lem:localweight}, or \(C=0\) if \(U=0\).
Every obstruction satisfies \(|Z[k]|\le |R_{r+mk}|+1\) eventually.
Its projective height slope is therefore at most \(D_0+C\).
Translate the joint support of its reduced pair by the inverse of
one of its characters to make it contain \(1\). The height slope
\(\sum_{v\in S}\max\{w_v(A),w_v(B)\}\) is unchanged. At each place
its local summand is now nonnegative and bounds
\(\log^+|\gamma|_v\) for every character \(\gamma\) in the
translated support. Thus \(h(\gamma)\le D_0+C\).
Northcott's finiteness theorem \cite[Lemma~2]{Northcott1950}
gives finitely many possible characters in \(K\), hence finitely
many possible supports and slopes within the bound. The \(T\)-degrees
do not enter \eqref{eq:alg-D}.

Two distinct obstructions of the same slope \(D\) have distinct
rational values for all sufficiently large \(k\), by
Lemma~\ref{lem:alg-orbit}. Rational separation then bounds the
distance of these values below by
\(\exp(-2Dk+o(k))\). On the other hand, the triangle inequality and
the strict obstruction inequalities bound their distance above by
\(O(\exp(-(2D+\eta)k))\) for some \(\eta>0\). This also applies if
one function is exact. It proves uniqueness at each obstruction
slope. In particular a nonempty obstruction set has a minimum
attained by a unique function.

For a function of slope \(D\), its reduced denominator exceeds
\(Q^{r+mk}\) eventually if \(D>m\log Q\), and is smaller eventually
if \(D<m\log Q\). Below the stated cutoff the finite list precludes
infinitely many violations at any fixed error tolerance. Above
the cutoff the minimizing function on one residue gives violations.
At a finite cutoff every higher-slope obstruction is eventually
excluded. The minimizing function on each minimizing residue gives
violations for a fixed tolerance if its reduced denominator fits
infinitely often, and otherwise is eventually excluded. These
observations prove \eqref{eq:alg-endpoint}, including the factor
\(Q_*^r\).
\end{proof}

\begin{corollary}\label{cor:alg-rational}
On a residue where \(R_{r+mk}\) takes rational values infinitely
often, \(U/V\) is invariant, \(d_r\le D(U/V)\), and
\eqref{eq:alg-admissibility} with the range replaced by
\(1\le q<\operatorname{den}(R_{r+mk})\) holds if and only if
\(d_r=D(U/V)\).
\end{corollary}

\begin{proof}
Invariance follows from Lemma~\ref{lem:alg-orbit}, and the exact
function is an obstruction. Write \(D_R=D(U/V)\).
If \(d_r<D_R\), the minimizing obstruction has reduced
denominator strictly below the exact denominator eventually and
gives violations. If \(d_r=D_R\), any distinct fixed invariant
function \(Z\) has, by rational separation,
\[
 2D(Z)+e(Z)\ge D(Z)-D_R.
\]
Consequently every distinct obstruction would have \(D(Z)<D_R\),
which is impossible. The range under consideration is eventually
contained in any fixed exponential ceiling with \(m\log Q>D_R\).
Lemma~\ref{lem:alg-finite} now allows only the exact function at
infinitely many violations, and its reduced denominator exceeds
every denominator in this range.
\end{proof}

\begin{remark}[Zero slope]\label{rem:alg-zero}
The identity
\[
 D(Z)=\sum_{v\in S}w_v(B)
       +\sum_{v\in S_f}\max\{0,w_v(A)-w_v(B)\}
\]
shows that \(D=0\) precisely when the denominator has one character
and, after division by it, every numerator character is an
algebraic integer. Indeed the first sum is nonnegative by the
product formula and bounds below the height of the ratio of any
two denominator characters. Distinct characters have a ratio of
positive height, by torsion-freeness and Northcott. After division
by the denominator character the remaining finite-place conditions
are exactly integrality of the numerator characters. The polynomial
factor in the denominator may be retained. A rational polynomial
and a constant clear the \(K(T)\)-coefficient denominators of
the resulting sum, so multiplying its rational values by that
polynomial gives integers at all sufficiently large \(k\).
In particular \(\operatorname{den}(Z[k])=O(k^C)\) for some \(C\).
The endpoint at \(Q=1\) requires the actual denominator. For example,
\(R_n=1/n+\sqrt{2}\,2^{-n}\) has zero-slope obstruction \(1/n\),
whose denominator is \(n\). Thus the cutoff is \(1\) and \(Q=1\)
is admissible.
\end{remark}

\begin{remark}
If \(U/V\) differs from its complex conjugate on a residue, let
\(\tau\in G\) be ordinary complex conjugation. For rational \(Z[k]\),
\[
 2|R_{r+mk}-Z[k]|\ge |R_{r+mk}-\overline{R_{r+mk}}|.
\]
The right side has exponential rate
\(\rho(U\tau V-\tau U\,V)-\rho(V)-\rho(\tau V)\) by
Lemma~\ref{lem:localweight}. Hence \(2D(Z)\) for every obstruction
is strictly smaller than the negative of that rate. The support
bound in the proof of Theorem~\ref{thm:algcut} then applies to all
obstructions on such a residue.
\end{remark}

\subsection{Fibonacci and cubic trace quotients}

\begin{proposition}\label{prop:fib-scalar}
For \(R_n=F_{n+1}/F_n\), with \(F_0=0,F_1=1\) and
\(F_{n+2}=F_{n+1}+F_n\), the cutoff is
\(\varphi=(1+\sqrt5)/2\) and the endpoint is inadmissible.
The bound \eqref{eq:alg-admissibility} holds throughout
\(1\le q<F_n\) at sufficiently large \(n\).
\end{proposition}

\begin{proof}
Consecutive Fibonacci numbers are coprime by the Euclidean
algorithm, so \(F_n\) is the exact denominator. Binet's formula
with \(\psi=-1/\varphi\) gives \(F_n\sim\varphi^n/\sqrt5\) and
\(|R_n-\varphi|\ll\varphi^{-2n}\).
The integral polynomial \(X^2-X-1\) does not vanish at \(p/q\),
so \(|(p/q-\varphi)(p/q-\psi)|\ge q^{-2}\). For
\(|p/q-\varphi|\le1\) this gives a constant \(c>0\) with
\(q^2|p/q-\varphi|\ge c\), and the inequality also holds outside
that neighborhood after decreasing \(c\).
For \(q\le c_1\varphi^n\), with \(c_1>0\) sufficiently small,
the triangle inequality gives \(q^2|R_n-p/q|\ge c/2\).
For \(c_1\varphi^n<q<F_n\), rational separation gives
\(q^2|R_n-p/q|\ge q/F_n\gg1\). This proves the full
strict-denominator bound and admissibility for \(Q<\varphi\).
The exact rational value fits at \(Q=\varphi\) for all sufficiently
large \(n\), since \(F_n<\varphi^n\).
\end{proof}

\begin{proposition}[The cubic Pisot trace ratio]\label{prop:cubic-scalar}
Let \(\theta>1\) be the real root of \(X^3-X-1\), with other
roots \(\beta,\gamma=\overline\beta\), and put
\(T_n=\theta^n+\beta^n+\gamma^n\) and \(R_n=T_{n+1}/T_n\).
Then the cutoff is \(\theta\), its endpoint is inadmissible,
and the bound \eqref{eq:alg-admissibility} holds throughout
\(1\le q<\operatorname{den}(R_n)\) at sufficiently large \(n\).
Moreover
\[
 T_n=\theta^n\bigl(1+O(\theta^{-3n/2})\bigr),\qquad
 R_n-\theta=O(\theta^{-3n/2}),
 \qquad \beta/\gamma\text{ is not a root of unity}.
\]
\end{proposition}

\begin{proof}
The cubic is irreducible by the rational root test. Its discriminant
is \(-23\), so its splitting field \(K\) has group \(S_3\) and
it has one real root and a conjugate pair. The real root is greater
than \(1\), the root product is \(1\), and all roots are algebraic
units. Hence \(|\beta|=|\gamma|=\theta^{-1/2}\). Traces are
integers. The estimates for \(T_n\) and \(R_n-\theta\) follow by
separating the real-root term in \(T_n\) and in \(T_{n+1}-\theta T_n\).

Write \(\alpha_1=\theta,\alpha_2=\beta,\alpha_3=\gamma\).
There are three complex places in \(K\), indexed so that at place
\(i\) the root \(\alpha_i\) has ordinary modulus \(\theta\)
and the other two have modulus \(\theta^{-1/2}\).
With \(\bar a=(a_1+a_2+a_3)/3\), the normalized weights are
\[
 \log\left|\prod_{j=1}^3\alpha_j^{a_j}\right|_i
 =\tfrac12(\log\theta)(a_i-\bar a).
\]
If a product is torsion all three weights vanish, so all \(a_i\)
are equal and the product is \(1\). This proves the assertion on
\(\beta/\gamma\) and identifies the character lattice with
\[
 L=\{a-\bar a(1,1,1):a\in\mathbb Z^3\}.
\]
Equivalently, \(L\) consists of the centered triples whose pairwise
differences are integers. We can take \(m=1\).

Consider a nonzero invariant function \(Z=A/B\) with the pair
reduced in \(K[T][\Lambda]\). For each \(\sigma\in S_3\),
coprimality and invariance give
\((\sigma A,\sigma B)=c_\sigma[\delta_\sigma](A,B)\), where
\(c_\sigma\in K^\times\) and \(\delta_\sigma\in L\) in additive
lattice notation. If \(b\) is the barycenter of the denominator
support, then \(\delta_\sigma=\sigma b-b\).
For the transposition of coordinates \(i,j\), the third coordinate
of \(\sigma b-b\) is zero. Since this triple belongs to \(L\),
its \(i\)-th coordinate \(b_j-b_i\) is an integer. Hence all
differences of the coordinates of \(b\) are integers, and the
coordinates sum to zero. Thus \(b\in L\).
Translate the common pair by the character with vector \(-b\).
The supports are now permutation invariant, and the Galois action
on the pair is multiplication by \(c_\sigma\) alone. These constants
satisfy \(c_{\sigma\tau}=c_\sigma\sigma(c_\tau)\). By multiplicative
Hilbert 90 \cite[Chapter~II, Proposition~1.22]{Hilbert90Milne}, one can rescale the pair by a
constant so that both polynomials are invariant. The same
normalization applies with \(K[T]\)-coefficients, since the common
units are constants times characters.

Let \(M\) be the largest first coordinate of the translated
denominator support. The maxima in the other two coordinates
equal \(M\), and \(M\in\frac13\mathbb Z_{\ge0}\) because \(L\) is
centered. With all characters units, \eqref{eq:alg-D} becomes
\(D(Z)=\frac32(\log\theta)M\). Thus the only slopes below
\(\log\theta\) are \(0\) and \(\tfrac12\log\theta\).

An obstruction converges exponentially to \(\theta\), by the
error estimate for \(R_n-\theta\) and Lemma~\ref{lem:alg-slopes}.
In slope zero the normalized denominator support is \(\{0\}\),
since every coordinate of each support point is at most zero.
Boundedness of \(Z[k]\) and Lemma~\ref{lem:localweight} force the
largest first coordinate of the numerator support to be at most
zero. Permutation invariance gives this bound in every coordinate.
Thus \(A,B\in\mathbb Q[T]\). A bounded rational function in
\(\mathbb Q(T)\) has a rational limit as \(k\to\infty\), which
excludes convergence to \(\theta\). The zero function is also
excluded by that convergence.

For the other slope, \(M=1/3\). A centered triple in \(L\) with
maximum coordinate \(\le1/3\) is either zero or a permutation of
\((-2/3,1/3,1/3)\). To see this, a nonzero triple has maximum
\(1/3\), and every coordinate differs from \(1/3\) by a nonnegative
integer with the three differences summing to \(1\).
The three nonzero triples correspond to the inverse-root
characters, and all three occur in \(B\) by permutation
invariance. Boundedness and the local weight lemma give the same
coordinate bound for \(A\), first at coordinate \(1\) and then at
the others by permutation invariance.

At the specified embedding the characters \([\beta^{-1}]\)
and \([\gamma^{-1}]\) are precisely the possible characters of
largest weight \(\tfrac12\log\theta\). Exponential convergence
forces the coefficients of both to vanish in \(A-\theta B\).
Indeed, if either survived, Lemma~\ref{lem:localweight} would
give exponential rate zero for \((A-\theta B)[k]/B[k]\), even
with tied weights. Let \(a_2(T),b_2(T)\) be the coefficients
of \([\beta^{-1}]\). Then \(a_2=\theta b_2\) and \(b_2\ne0\).
Both coefficients are fixed by the transposition that fixes
\(\beta\) and exchanges \(\theta\) with \(\gamma\).
Applying that transposition gives \(a_2=\gamma b_2\), a
contradiction. This excludes every obstruction below \(\log\theta\).

Since the exact rational value has denominator at most
\(|T_n|\), its slope is at most \(\log\theta\), and it is an
obstruction. Theorem~\ref{thm:algcut} and
Corollary~\ref{cor:alg-rational} give the cutoff and the
strict-denominator bound. Finally write
\(\beta=\theta^{-1/2}e^{i\omega}\). If \(\omega/(2\pi)\) were
rational then \(\beta/\gamma=e^{2i\omega}\) would be torsion.
By density of the irrational rotation, \(\cos(n\omega)<0\)
infinitely often. At those sufficiently large levels
\(0<T_n=\theta^n+2\operatorname{Re}(\beta^n)<\theta^n\), so the exact
rational value fits the ceiling at \(Q=\theta\) and violates
\eqref{eq:alg-admissibility}.
\end{proof}

\begin{remark}\label{rem:cubic-localweight}
For the cubic roots of Proposition~\ref{prop:cubic-scalar}, any
nonzero Laurent polynomial \(L_0\) over the splitting field satisfies
\(\log|L_0((\beta/\gamma)^n)|=o(n)\). This is
Lemma~\ref{lem:localweight} on the cyclic group generated by
\(\beta/\gamma\): its distinct powers are distinct characters
and all have ordinary modulus \(1\) at the prescribed embedding.
\end{remark}

\section{Simultaneous dynamical recurrence on the diagonal}\label{sec:rec}
Let \(b_1,\ldots,b_k\ge2\) be pairwise coprime integers and \(N_i(n)=b_i^n-1\). For \(0<\psi_i(n)\le1\) and arbitrary \(y_{i,n}\in\mathbb R\), let \(E\) consist of those \(x\in\mathbb R\) for which
\begin{equation}\label{eq:rectests}
                   \|N_i(n)x-y_{i,n}\|<\psi_i(n)\qquad (1\le i\le k)
\end{equation}
at infinitely many common indices. At zero phase these are returns of \(b_i^nx\bmod1\) to the starting point \(x\bmod1\). Allowing the starting coordinates to vary separately gives an ambient torus problem; see, for example, \cite{ambient}. Here every coordinate has the same starting parameter.
\begin{theorem}\label{thm:rec}
Put \(\kappa_i=\log b_i\). At \(g=(g_1,\ldots,g_k)\in[0,\infty)^k\) define \(t_i=\kappa_i+g_i\) and
\begin{equation}\label{eq:recformula}
 d(g)=\left[1-\max_{\varnothing\ne S\subseteq\{1,\ldots,k\}}
        \frac{\sum_{i\in S}g_i}{\max_{i\in S}t_i}\right]_+.
\end{equation}
Set \(d(g)=0\) if any coordinate is infinite. For every nonempty open interval \(J\),
\[
             \dim_H(J\cap E)=\sup_{g\in\mathcal C}d(g),
\]
where \(\mathcal C\) is the joint cluster set in \([0,\infty]^k\) of the vectors \((-n^{-1}\log\psi_i(n))_{i=1}^k\). Each symmetric test may also be replaced by a one-sided open interval of fractional length \(\psi_i(n)\). More generally one may take \(N_i(n)=F_i(b_i^n)\) for nonconstant \(F_i\in\mathbb Z[X_i]\) with positive leading coefficients, using \(\kappa_i=(\deg F_i)\log b_i\).
\end{theorem}

\begin{lemma}\label{lem:recdual}
In the polynomial case of Theorem~\ref{thm:rec}, order the axes arbitrarily. If \(u,g_i\ge0\) and
\begin{equation}\label{eq:recdens}
               u+\sum_{i<j}g_i<\kappa_j\qquad(1\le j\le k),
\end{equation}
then for each fixed \(C>0\) and all sufficiently large \(n\) there is no \(0\ne h\in\mathbb Z^k\) with \(|h_i|\le C e^{g_i n}\) and \(|\sum_iN_i(n)h_i|\le C e^{un}\).
\end{lemma}
\begin{proof}
Use \(X_{i,n}=b_i^n\) and monomial weights \(w(X_i)=\log b_i\). From an infinite sequence of counterexamples choose a subspace \(W\) over \(\mathbb Q(X_1,\ldots,X_k)\) of minimal rank \(r\) whose specializations contain infinitely many of the vectors. On this subsequence let \(P\) be a primitive polynomial Pl\"ucker vector of \(W\), so \(P_I\) is indexed by \(r\)-element coordinate sets. The identity chart in Theorem~\ref{thm:descent} gives \(w(P_I)\le\sum_{i\in I}g_i\). For \(|J|=r-1\), contract \(P\) with the row \(L=(F_i(X_i))_i\), giving
\[
                     L_J=\sum_{i\notin J}\pm F_i(X_i)P_{J\cup i}.
\]
Take any \(l\notin J\). Replace coordinate \(h_l\) by \(Lh/F_l\), leaving the other coordinates fixed. The \(J\cup\{l\}\) minor is now \(\pm L_J/F_l\) and the new coordinate has cap exponent \(u-\kappa_l\). The same descent theorem gives \(w(L_J)\le u+\sum_{j\in J}g_j\), since a strict violation in either chart would decrease the rank.

For \(I,J\subseteq[k-1]\) both bounds are strictly less than \(\kappa_k\), by \eqref{eq:recdens}. Hence those polynomials have \(X_k\)-degree below \(\deg F_k\). In \(L_J\) every term with index below \(k\) retains this degree bound. The remaining term \(F_k(X_k)P_{J\cup k}\) must vanish. This eliminates every minor containing \(k\). Descending through the remaining indices eliminates all minors by the same argument, contrary to \(r\ge1\).
\end{proof}

\subsection{Counting and dimension}
First take fixed power radii \(e^{-g_i n}\). Fixed positive factors in the widths do not change any estimate below. Order the indices by \(t_1\le\cdots\le t_k\), and set
\[
           D_j=\kappa_j-\sum_{i<j} g_i,\qquad D=D_k,\qquad m_0=\min_jD_j .
\]
At a terminal \(t_j\), the largest numerator in \eqref{eq:recformula} includes every coordinate ending by that terminal. Among indices with the same \(t_j\), the last index minimizes \(D_j/t_j\). Thus \mbox{\(d(g)=[\min_j D_j/t_j]_+\)}. Assume first that \(m_0>0\). On each nonempty subsystem \(S\) keep the induced order, and denote its corresponding minimum, last count exponent and last terminal by \(m_S,D_S,t_S\). In particular \(m_S>0\).

For \(0\le u<m_S\) consider in \(\mathbb R^S\) the symmetric body
\[
 K_S(u)=[-e^{-un},e^{-un}](N_i(n))_{i\in S}
                  +\prod_{i\in S}[-e^{-g_i n},e^{-g_i n}].
\]
Its width in an integral direction \(h\) is
\(2(e^{-un}|\sum_{i\in S}N_i h_i|+\sum_{i\in S}e^{-g_i n}|h_i|)\).
The strict inequalities in Lemma~\ref{lem:recdual} persist after increasing \(u\) and each \(g_i\) by a sufficiently small \(\eta>0\). Thus the minimum integral width of \(e^{-\eta n}K_S(u)\) tends to infinity. By the flatness theorem \cite{KL}, every translate of this smaller body contains a lattice point for all large \(n\). Its lattice translates cover the space, so it contains a measurable fundamental domain: partition the unit cube according to the first covering translate in an enumeration and translate each piece back.

Fix positive transverse constants and an interval of parameter length \(v\ge e^{-un}\) in a bounded range. Sweep the transverse box of those constant multiples of \(e^{-g_i n}\) along the image of the interval under \(x\mapsto(N_i(n)x)_{i\in S}\), translating by the phases. Its volume is \(\asymp v e^{D_S n}\): the box volume is augmented by the sum of the coordinate displacements times their transverse face volumes, and the largest swept term has terminal \(t_S\). Here \(u<m_S\le t_S\). The lattice-point count is likewise \(\asymp v e^{D_S n}\), uniformly in position and phases. For the upper bound attach fundamental cells and enlarge the segment and box by fixed factors; for the lower bound shrink both by fixed factors and cover the shrunken body by cells whose lattice points belong to the original one. The exponentially smaller cell fits the margins in both comparisons.

Apply this count to the first \(j\) coordinates on a bounded parameter interval, using \(u=0\) and, if necessary, enlarging that interval to a unit interval. There are \(O(e^{D_jn})\) integer tuples witnessing a hit and each contributes an interval of length \(O(e^{-t_jn})\). Summing the best prefix cover over the levels proves the upper dimension bound for fixed positive \(m_0\).

For the lower bound, use any prescribed infinite sequence of levels. In a parent interval of length \(R\), count the witnesses in a tube over its middle half, with transverse widths reduced by a sufficiently small fixed factor. Center one-sided targets at their midpoints. Each witness yields a common interior interval of length \(c e^{-t_kn}\) about its root on an axis ending at \(t_k\): the displacement from a witnessing parameter is at most a small multiple of \(e^{-t_kn}\), which fits inside the other physical windows. The selected centers at distance \(O(e^{-t_k n})\) from any one point have bounded multiplicity, since every numerator varies there by \(O(1)\). For large \(n\) one may therefore keep \(\asymp R e^{Dn}\) children in the parent with gaps at least their lengths.

We bound the number of these children meeting an interval of length \(e^{-zn}\) with \(0\le z<t_k\). Enlarge the test interval by \(O(e^{-t_kn})\). The axes ending by \(z\) each have \(O(1)\) possible numerators. Put \(S=\{i:t_i>z\}\), \(G=\sum_{i\notin S}g_i\); thus \(D_S=D+G\). In the remaining coordinates apply the tube upper bound after enlarging the length to \(e^{-un}\), with \(u\le\min(z,m_S)\) and strictly below \(m_S\). For every \(\delta>0\) this gives the bound
\begin{equation}\label{eq:reclocal}
 O_\delta\big(e^{(D-E(z)+\delta)n}\big),\qquad
 E(z)=\min\big\{z-G,\ \min_{t_j>z}D_j\big\}.
\end{equation}
Indeed \(m_S=G+\min_{t_j>z}D_j\). A fixed mesh of choices for \(u\) from below, including zero and leaving a small margin at \(m_S\), makes the estimate uniform in \(z\). The first entry \(z-G\) in \eqref{eq:reclocal} is at least \(d(g)z\). At the last terminal by \(z\) its value is the corresponding \(D_j\); until the next terminal it grows with slope one. Below the first terminal \(G=0\). The other entries satisfy \(D_j\ge d(g)t_j\). For coarse lengths \(v\ge e^{-z_0n}\), with \(0<z_0<m_0\) fixed, the full tube count gives the stronger bound \(O(v e^{Dn})\).

Fix \(0<s<d(g)\). Choose sparse generations from the prescribed levels, dividing each parent mass \(P\) equally among its children. Use an initial fixed normalization and the induction hypothesis \(P\le C_0R^s\). In lengths \(e^{-z_0n}\le v\le R\) the mass meeting a test interval from this parent is \(O((P/R)v)=O(v^s)\). For \(e^{-t_kn}\le v< e^{-z_0n}\), \eqref{eq:reclocal} bounds it by \(O_\delta((P/R)e^{\delta n}v^{d(g)})\). Take \(\delta<z_0(d(g)-s)/2\) and then \(n\) sufficiently large for the finite parent generation. At lengths from \(c e^{-t_kn}\) to \(e^{-t_kn}\), bounded multiplicity applies. Each child mass is \(O((P/R)e^{-Dn})\le C_0(c e^{-t_kn})^s\) at sufficiently late levels, since \(D\ge d(g)t_k>s t_k\). This preserves the induction. An interval no longer than the previous parent length meets a bounded number of parents. The resulting measure gives the claimed lower bound in every open interval by the mass distribution principle.

To conclude the proof of Theorem~\ref{thm:rec}, note that \eqref{eq:recformula} is continuous for finite rates. So is
\(m_0=\min_{S\ne\varnothing}(\max_{i\in S}t_i-\sum_{i\in S}g_i)\).
When \(m_0\le0\), move along the segment from \(0\) to \(g\) to the first vanishing of \(m_0\) and use the upper covers just inside the positive range. They cover the prescribed tests and have dimensions tending to zero. For arbitrary radii any finite joint limit gives a lower bound along a subsequence by increasing all limiting exponents slightly. For the upper bound fix \(s>\sup_{\mathcal C}d\). Near every finite cluster vector, enlarge the windows to slightly weaker powers so that one of the upper covers has exponentially decreasing \(s\)-cost. Coordinates of rate zero can use a full-width test. If the vector is on or beyond the zero-dimensional boundary, use the preceding scaling argument and then weaken positive rates slightly to cover a neighborhood. Near an infinite coordinate the single-grid cover has exponentially decreasing cost for any \(s>0\). A finite subcover of the compact cluster set contains all sufficiently late levels; summing in \(n\) finishes the proof.

\subsection{Three bases and the effect of shared primes}
For \(2,3,5\) and common \(g\), Theorem~\ref{thm:rec} gives
\[
 \left[\min\left\{\frac{\log 2}{\log 2+g},
                  \frac{\log 3-g}{\log 3+g},
                  \frac{\log 5-2g}{\log 5+g}\right\}\right]_+ .
\]
At \(g=\log 2\) the last term gives \(\log(5/4)/\log 10\). At \(g=\tfrac12\log 5\) the dimension is zero although each of the three pairs still has positive dimension.

\begin{proposition}\label{prop:shared}
Let \(1<c<a\) be integers and \(b=ac\). Put \(\kappa=\log c,\ \alpha=\log a,\ \beta=\log b\). If \(\kappa<g<\alpha\), the two recurrence tests for \(a,b\), with common radius \(e^{-gn}\), have local dimension
\[
                   \min\left\{\frac{\kappa}{g},\frac{\beta-g}{\beta+g}\right\}
\]
for all phases, including one-sided windows.
\end{proposition}
\begin{proof}
Write \(N=a^n-1,\ M=b^n-1,\ D'=c^n-1,\ \delta=e^{-gn}\). The identity \(M=c^nN+D'\) forces every hit to lie within \(O(\delta)\) of the roots of \(D'x=y_{2,n}-c^n y_{1,n}\bmod1\). There are \(O(e^{\kappa n})\) such clusters on a fixed compact interval. Each contains \(O(N\delta+1)\) first-grid centers, each supporting \(O(1)\) fine intervals of length \(O(\delta/M)\), since \(M\delta/N\to0\). These two covers give the upper bounds.

Inside a small concentric subcluster of radius \(c_0\delta\) take first-grid centers \(x_i\) for the midpoint of the first window, using midpoint phases in the cluster equation. Their displacement to a nearest fine center for the second midpoint is \(O(c_0D'\delta/M)=O(c_0\delta/N)\). Hence each supports a common interior interval of length comparable to \(\delta/M\). Inside the middle half of a fixed parent of length \(R\), at all sufficiently late levels there are \(\asymp R e^{\kappa n}\) full subclusters, each with \(\asymp e^{(\alpha-g)n}\) separated children at mesh comparable to \(e^{-\alpha n}\). Give equal mass to the children and write \(P\) for the parent mass. Between the child and parent scales, an interval of length \(v\) has mass from this parent at most a constant times \(P/R\) multiplied respectively by
\[
 \begin{cases}
 v,&v\ge e^{-\kappa n},\\
 e^{-\kappa n},& e^{-gn}\le v<e^{-\kappa n},\\
 v e^{(g-\kappa)n},&e^{-\alpha n}\le v<e^{-gn},\\
 e^{-(\beta-g)n},&v<e^{-\alpha n}.
 \end{cases}
\]
The first two estimates count clusters and the last two count children within a cluster. For \(s\) below both asserted bounds the second and fourth estimates have strict slack relative to \(v^s\); for the third use \(v^{1-s}\le e^{-(1-s)gn}\). The sparse-generation argument above gives a measure with the desired \(s\)-bound.
\end{proof}
For \((a,b)=(81,162)\) at radius \(8^{-n}\), Proposition~\ref{prop:shared} gives \(1/3\). The coprime expression is strictly larger: its second fraction exceeds \(1/3\) because \(162>8^2\), and its first does because \(81^2>8\).

For bases \((2,3,6)\) take fractional rates \(g_1=\log6,\ g_2=\log(3/2),\ g_3=g_2/2\). Writing \(N_b=b^n-1\), the identity \(N_6=3^nN_2+N_3\) implies at each zero-phase hit of the first two tests that
\[
                \|N_6x\|\le 2^{-n}+(2/3)^n=o((2/3)^{n/2}).
\]
The third zero-phase test is eventually automatic. Equation~\eqref{eq:recformula} on the first two axes gives \(\log(4/3)/\log12\). Shifting only the third target by \(1/2\) makes every sufficiently late intersection empty. All three pairs are multiplicatively independent and have subexponential recurrence gcd by \cite{BCZ}. For general shared-prime bases with overlapping scales, a dimension formula accounting for these joint phase constraints remains open.

\section{Polynomial values and algebraic approximants}
\label{sec:higher-degree}

For \(0\ne P=\sum_i a_i Z^i\in\mathbb Z[Z]\), write
\(H(P)=\max_i|a_i|\). For an algebraic number \(\alpha\), let
\(P_\alpha\) be its primitive integral minimal polynomial and set
\(H(\alpha)=H(P_\alpha)\). Thus a reduced rational \(p/q\) with
\(q>0\) has height \(\max(|p|,q)\). For an integer \(d\ge1\) and
\(Q\ge1\), the two bounds considered here are
\begin{align}
  |P(R_n)|H(P)^d &\ge \exp(-\epsilon n)
  &&(0\ne P,\ \deg P\le d,\ H(P)\le Q^n), \label{eq:higher-poly}\\
  |R_n-\alpha|H(\alpha)^{d+1} &\ge \exp(-\epsilon n)
  &&(\deg\alpha\le d,\ H(\alpha)\le Q^n). \label{eq:higher-root}
\end{align}
Each is required for every \(\epsilon>0\) and all sufficiently large
\(n\), uniformly over the indicated approximants. Exact hits are
included unless explicitly excluded. A cutoff of \(1\) can therefore
occur with failure at \(Q=1\). Theorem~\ref{thm:algcut} treats
rational approximants with a denominator ceiling.

\begin{lemma}[Evaluation near a root]\label{lem:higher-transfer}
Let \(\alpha\) have degree \(e\le d\). If \(|z-\alpha|\le1\), then
\[
 |P_\alpha(z)|\ll_d
 H(\alpha)(1+|z|)^{e-2}|z-\alpha|.
\]
\end{lemma}

\begin{proof}
If \(a\) is the leading coefficient and the roots are
\(\alpha=\alpha_1,\ldots,\alpha_e\), Jensen's formula and the
quadratic mean on the unit circle give
\[
 |a|\prod_j\max(1,|\alpha_j|)
 \le \Big(\sum_i |[Z^i]P_\alpha|^2\Big)^{1/2}
 \le \sqrt{e+1}\,H(\alpha).
\]
In the factorization of \(P_\alpha(z)\), use
\(|z-\alpha_j|\le(1+|z|)\max(1,|\alpha_j|)\) for \(j\ge2\).
The remaining factor \(\max(1,|\alpha|)\) from the displayed
product is at least \((1+|z|)/3\), since \(|z-\alpha|\le1\).
\end{proof}

\subsection{Positive integer roots}

Let \(k\ge1\), let \(s_1,\ldots,s_k>1\) be pairwise coprime integers,
and put \(X_n=(s_1^n,\ldots,s_k^n)\) and
\(F=\mathbb Q(X_1,\ldots,X_k)\). Quotients of power sums with
positive integer roots and rational coefficients have this form
by choosing the primes in the roots as generators. For a
nonzero polynomial \(G\) in the \(X_i\), put
\[
 w(G)=\max_{\boldsymbol a\in\operatorname{supp}G}
       \sum_{i=1}^k a_i\log s_i .
\]
Set \(w(0)=-\infty\) and extend \(w\) to \(F^\times\) by subtraction.
Distinct monomials have distinct weights. In particular,
\(\log |G(X_n)|=nw(G)+O_G(1)\) for every \(G\in F^\times\)
and all sufficiently large \(n\).

Fix \(R=U/V\in F^\times\), with \(U,V\) coprime polynomials,
and let \(R_n=R(X_n)\). For \(\eta\in\mathbb R\), call \(Q\ge1\)
\emph{admissible} if for every \(\epsilon>0\), eventually
\begin{equation}\label{eq:higher-scaled}
 |P(R_n)|H(P)^d\ge \exp((\eta-\epsilon)n)
 \qquad(0\ne P,\ \deg P\le d,\ H(P)\le Q^n).
\end{equation}
If \(R=0\), the polynomial \(Z\) obstructs admissibility for
all \(\eta\), already at \(Q=1\). For \(R\ne0\) let
\(L=(1,R,\ldots,R^d)\). For each subspace
\(0\ne W\subset F^{d+1}\) of dimension \(l\), take a Pl\"ucker
vector \(\Pi\in\bigwedge^l F^{d+1}\) of \(W\) with polynomial
coordinates having no common nonconstant factor. It is unique
up to a nonzero rational factor. Define
\begin{equation}\label{eq:higher-DE}
 D(W)=\max_{|I|=l} w(\Pi_I),\qquad
 E(W)=\max_{|J|=l-1} w((\iota_L\Pi)_J),
\end{equation}
where \(\iota_L\) denotes contraction by the covector \(L\).
In particular \(E(W)=-\infty\) when \(L\) annihilates \(W\).

\begin{theorem}[Sharp polynomial cutoff]\label{thm:polycut}
For \eqref{eq:higher-scaled}, define
\begin{equation}\label{eq:higher-min}
 g_*=\min_{0\ne W\subset F^{d+1}}
 \left\{\frac{D(W)}{l}:
   E(W)+(d-l+1)\frac{D(W)}{l}<\eta,\quad l=\dim W\right\}.
\end{equation}
The minimum is attained. Every \(1\le Q<\exp(g_*)\) is
admissible and every \(Q>\exp(g_*)\) fails. At \(Q=\exp(g_*)\),
admissibility is equivalent to the following condition for every
minimizing \(W\). For all sufficiently large \(n\), the lattice
\(W(X_n)\cap\mathbb Z^{d+1}\) has no nonzero vector of sup norm
at most \(\exp(g_* n)\).
\end{theorem}

\begin{proof}
Write \(W_n=W(X_n)\). The ranks of the fixed rational-function
matrices used to define \(W_n\) stabilize because their nonzero
minors are eventually nonzero upon evaluation. Clear fixed
rational denominators in \(\Pi\). The basis wedge of
\(W_n\cap\mathbb Z^{d+1}\) is, up to sign, \(\Pi(X_n)\)
divided by the gcd of its coordinates. Indeed the lattice is
saturated, so an integral basis extends to a basis of
\(\mathbb Z^{d+1}\), and its wedge is primitive.

The gcd is \(\exp(o(n))\) by Lemma~\ref{lem:evaluated-gcd}.

For real \(g\) and \(\delta>0\), send the lattice into
\(\mathbb R^{d+2}\) by
\begin{equation}\label{eq:higher-map}
 {\bf a}\longmapsto
 \big((e^{-gn}a_i)_{i=0}^d,\ 
             e^{(dg+\delta-\eta)n}L(X_n){\bf a}\big).
\end{equation}
Its covolume in its \(l\)-dimensional real span is the
Euclidean norm of the image of the basis wedge. The minors
using only the coefficient rows have maximum weight
\(D-lg\) after scaling. The minors using the evaluation row
have maximum weight \(E+(d-l+1)g+\delta-\eta\).
Consequently the logarithmic covolume divided by \(n\)
converges to
\begin{equation}\label{eq:higher-covol}
 \max\{D-lg,\ E+(d-l+1)g+\delta-\eta\}.
\end{equation}

Suppose \eqref{eq:higher-scaled} fails for some \(\epsilon>0\).
Pass to an infinite sequence of violating coefficient vectors
\({\bf a}_n\) such that \(\log H(P)/n\to g\le\log Q\).
Fix \(0<\delta<\epsilon\). On this sequence the coefficient
coordinates have exponents at most \(g+o(1)\), and the
evaluation has exponent at most \(\eta-dg-\delta+o(1)\).
Start with \(W=F^{d+1}\). Theorem~\ref{thm:descent} in
complementary-minor form gives a proper subspace of the
current \(W\) on a further infinite sequence whenever
\eqref{eq:higher-covol} is positive, as follows. For \(D>lg\)
use its identity chart and the coefficient bounds. For a
positive term from a mixed minor with index set \(J\) on the
coefficient rows, choose \(j\notin J\). In the square chart
replace coordinate \(a_j\) by \(L{\bf a}/R^j\), retaining the
other coordinates. This is a permitted replacement with
pivot coefficient \(1\). The \((J\cup\{j\})\)-minor of
the transformed wedge is, up to sign,
\((\iota_L\Pi)_J/R^j\). The exponent bound for the
replacement coordinate is \(\eta-dg-\delta-w(R^j)+o(1)\),
so the same normalization cancels from the difference
between the minor weight and the sum of the coordinate
bounds. One may enlarge the bounds by a fixed small
positive amount in either chart, preserving the positive
difference and absorbing the \(o(1)\) errors.

At each step the dimension strictly decreases, and
\({\bf a}_n\ne0\) excludes the zero space. At termination
\[
 D\le lg,\qquad E+(d-l+1)g\le\eta-\delta .
\]
Since \(d-l+1\ge0\), this gives a subspace feasible in
\eqref{eq:higher-min} with slope at most \(g\). In particular
no violation is possible when \(\log Q<g_*\).

For the converse, take a feasible \(W\) with \(D/l<\log Q\).
Choose \(D/l<g<\log Q\) close enough to \(D/l\), and then
\(\delta>0\) small enough, so that both terms in
\eqref{eq:higher-covol} are strictly negative. The unit cube
of \(\mathbb R^{d+2}\) intersects the real image of \(W_n\)
in a centrally symmetric convex body containing its
\(l\)-dimensional Euclidean unit ball. Its volume thus
exceeds \(2^l\) times the lattice covolume for all large \(n\).
Minkowski's theorem supplies a nonzero lattice point in the
cube. The corresponding polynomial has \(H(P)\le e^{gn}\)
and \(|P(R_n)|\le e^{(\eta-dg-\delta)n}\). Its normalized
value is at most \(e^{(\eta-\delta)n}\), contradicting
admissibility for every \(0<\epsilon<\delta\).

The exact kernel of \(L\) is always feasible. Its dimension
is \(d\), and its Pl\"ucker coordinates are, up to signs,
\(V^d,UV^{d-1},\ldots,U^d\). They have no common nonconstant
factor, so its slope is \(\max\{w(U),w(V)\}\).
The monomials of bounded weight form a finite set and
\(1\le l\le d+1\). Hence there are only finitely many slopes
at or below this exact-kernel slope, and the minimum is attained.

Finally, if \(\log Q=g_*\) and a violation occurs infinitely
often, the descent gives a fixed feasible \(W\) with
\(g_*\le D/l\le g\le g_*\). Thus a minimizing \(W\) contains
lattice points within the endpoint ceiling at infinitely many
levels. Conversely, fix a minimizing \(W\) and an index
set \(I\) with \(w(\Pi_I)=D\). Cramer's rule expresses
\(L|_W\) as a combination of the coordinate functions
indexed by \(I\), with coefficients of the form
\[
 \pm\frac{(\iota_L\Pi)_{I\setminus\{j\}}}{\Pi_I}
 \qquad(j\in I).
\]
Each has weight at most \(E-D\), and if \(E=-\infty\)
the evaluation vanishes on \(W_n\). Any lattice point of
height \(H(P)\le e^{g_*n}\) in \(W_n\) therefore has
\[
 |P(R_n)|H(P)^d
 \ll_W e^{(E-D+(d+1)g_*)n}
 =e^{(E+(d-l+1)g_*)n}.
\]
The exponent is strictly less than \(\eta\), with the
estimate interpreted as zero when \(E=-\infty\).
Such points at infinitely many levels violate admissibility
at the endpoint. Eventuality in the lattice criterion is
for each fixed minimizing subspace.
\end{proof}

\begin{corollary}[Linear systems in degree two]\label{cor:higher-linear}
For \(d=2\) and \(\eta=0\), finitely many rational linear
systems determine the cutoff in Theorem~\ref{thm:polycut}.
\end{corollary}

\begin{proof}
For \(l=1,2,3\), every nonzero element of
\(\bigwedge^l F^3\) is decomposable. For \(l=2\), this follows,
for example, by identifying \(\bigwedge^2 F^3\) with
\((F^3)^*\) using a volume form and taking the wedge of a
basis of the kernel. Enumerate \(D\) among the monomial
weights with \(D/l\le\max\{w(U),w(V)\}\). Allow as unknowns
all coefficients of monomials of weight at most \(D\) in
each coordinate of \(\Pi\). Multiplying \(\iota_L\Pi\)
by \(V^2\) gives polynomial coordinates linear in these
unknowns. The required weight inequality is equivalent to
vanishing of all their coefficients at weights at least
\[
  2w(V)-(3-l)D/l.
\]
A solution has \(w(\Pi)=D\) precisely when one of its
weight-\(D\) coefficients is nonzero. Testing this is a
linear-algebra calculation on the solution space. All
weight comparisons reduce to comparisons of integer
powers of monomial bases.

If a feasible polynomial tuple has a nonconstant common
factor of weight \(u>0\), cancellation decreases \(D\) by
\(u\) and decreases \(E\) by \(u\) when it is finite. It
therefore preserves the strict feasibility inequality
and lowers the slope. Thus the least feasible slope from
the systems equals the minimum over primitive tuples.
The endpoint still uses the lattice condition of
Theorem~\ref{thm:polycut}.
\end{proof}

\begin{remark}
For \(d\ge2\), set \(\eta=(d-2)\max\{0,w(R)\}\).
Lemma~\ref{lem:higher-transfer} gives
\eqref{eq:higher-root} at every admissible \(Q\) for this
choice of \(\eta\). For \(|R_n-\alpha|\le1\) the loss
from the factor \((1+|R_n|)^{e-2}\) is \(O(e^{\eta n})\),
and any fixed multiplicative constant is absorbed by
applying \eqref{eq:higher-scaled} with a smaller tolerance.
Distances greater than \(1\) already satisfy the bound.
The lattice construction in Theorem~\ref{thm:polycut}
proves sharpness for polynomial values. A sharp
root-distance cutoff would additionally require
control of the roots of the constructed polynomials.
\end{remark}


\subsection{Powers of a rational number}

\begin{theorem}\label{thm:ratpower}
Let \(a>b\ge1\) be coprime and put \(t_n=(a/b)^n\). For
each integer \(e\ge1\) and every \(\epsilon>0\), eventually
\begin{equation}\label{eq:higher-ratpoly}
 |P(t_n)|H(P)^e\ge t_n^e e^{-\epsilon n}
 \qquad(P\in\mathbb Z[Z],\ \deg P=e,\ P(t_n)\ne0).
\end{equation}
For every algebraic \(\alpha\ne t_n\) of degree \(e\)
with \(|t_n-\alpha|\le1\), eventually
\begin{equation}\label{eq:higher-ratroot}
 |t_n-\alpha|H(\alpha)^{e+1}\ge t_n^2 e^{-\epsilon n}.
\end{equation}
Both estimates are uniform without a height ceiling.
For every fixed \(d\ge1\), the cutoffs in
\eqref{eq:higher-poly} and \eqref{eq:higher-root} are \(a\),
with failure at \(Q=a\) from exact hits. In particular,
the two cutoffs for \((3/2)^n\) are \(3\).
\end{theorem}

\begin{proof}
Write \(A=a^n\), \(B=b^n\) and \(t=A/B\).
If \(P(t)\ne0\), then \(|P(t)|\ge B^{-e}\).
Thus a violation of \eqref{eq:higher-ratpoly} satisfies
\(H(P)<Ae^{-\epsilon n/e}\). Suppose such violations occur
infinitely often. Among rational linear subspaces of the
coefficient space containing the violating polynomials on
infinite subsequences, choose one of least dimension.
Choose a reduced echelon basis \(f_1,\ldots,f_s\) of this
subspace by eliminating coefficients from degree \(e\)
downward. The leading degrees are distinct, the leading
coefficients are \(1\), and the coefficients at the other
basis pivots are zero. For \(P=\sum_i z_i f_i\),
the \(z_i\) are therefore integral and \(|z_i|\le H(P)\).
On a further infinite subsequence, all \(z_i\ne0\),
by minimality of the dimension. Let \(p\) be the index
with \(\deg f_p=e\). If \(s=1\), then
\(|P(t)|/t^e\gg |z_p|\ge1\) eventually, a contradiction.
We show that \(s\ge2\) also contradicts minimality.

Take an integer \(h\ge e\) and the coordinate vector
\((z_i t^j)_{1\le i\le s,\,0\le j\le h}\).
At the real place replace the coordinate forms of block
\(p\) for \(e\le j\le h\) by \(t^{j-e}P(t)\), viewed as
fixed rational linear forms in these coordinates. They
are triangular on block \(p\) with pivot \(1\), because
\(f_p\) is monic of degree \(e\). At the primes dividing
\(ab\), retain the coordinate forms. If \(I=h-e+1\),
the product over all chosen places is bounded by
\begin{align*}
 &\left(\prod_i|z_i|\right)^{h+1-I}
 \left(|P(t)|t^{-e}\prod_{i\ne p}|z_i|\right)^I\\
 &\qquad \le H(P)^{s(h+1-I)}e^{-\epsilon n I}.
\end{align*}
Here the product formula cancels the powers of \(t\),
the finite-place factors from \(z_i\) are at most \(1\),
and \(s-1\le e\). The vector is integral outside the
chosen places and has height at most exponential in \(n\)
for fixed \(h\). Taking \(\epsilon I>se\log a\) makes the
product exponentially small. The \(S\)-adic Subspace
Theorem \cite{Schlickewei} gives a nontrivial
relation on a further infinite subsequence,
\[
                 \sum_i z_i C_i(t)=0
\]
with fixed polynomials \(C_i\in\mathbb Z[Z]\), not all zero.
Here and below the unnormalized form-product conclusion
follows from the normalized theorem because the local
sup-norm product over the chosen places is at least the
projective height for vectors integral outside them.

If the formal combination \(\sum_i z_i C_i(Z)\) is
nonzero at one of these levels, remove its maximal
power of \(Z\). The rational root theorem then forces
its first nonzero coefficient to be a nonzero multiple
of \(A\). Its coefficients are \(O(H(P))=o(A)\), a
contradiction at all large levels. Thus the combination
vanishes formally on the subsequence. Coefficient
comparison imposes a fixed proper rational linear
restriction on \((z_i)\), contradicting minimality.
This proves \eqref{eq:higher-ratpoly}.

For \(\alpha\ne t\) with \(|t-\alpha|\le1\), apply
Lemma~\ref{lem:higher-transfer} and
\eqref{eq:higher-ratpoly} to \(P_\alpha\). Since \(t\to
\infty\), their combination yields
\(\lvert t-\alpha\rvert H(\alpha)^{e+1}
 \gg_e t^2 e^{-\epsilon n/2}\), which implies
\eqref{eq:higher-ratroot}. Distances greater than \(1\)
already satisfy \eqref{eq:higher-root}. Nonzero constant
polynomials satisfy \eqref{eq:higher-poly} directly.
Finitely many choices of \(e\le d\) give uniform bounds
in both forms with the exponents in
\eqref{eq:higher-poly}--\eqref{eq:higher-root}.

If \(P(t)=0\), the rational root theorem after removing
any factor \(Z\) gives \(H(P)\ge A\). The polynomial
\(BZ-A\) has height \(A\) and achieves an exact hit.
Likewise \(t\) is the exact algebraic approximant and
has height \(A\). This proves the cutoffs.
\end{proof}

\begin{remark}
For a positive rational less than \(1\), its powers
converge exponentially to \(0\). The height-one
approximants \(Z\) and \(0\) violate
\eqref{eq:higher-poly} and \eqref{eq:higher-root} even
with exact hits excluded. For the constant sequence
\(t_n=1\), exact hits occur at height one, and the
unrestricted estimates with exact hits excluded follow
from integrality and Lemma~\ref{lem:higher-transfer}.
Replacing the base by its negative preserves the
estimates by changing the signs of the polynomial
variable and the approximant on the odd levels.
\end{remark}

\begin{corollary}[Optimal nonexact exponents for rational powers]
\label{cor:higher-rat-exponents}
Under the hypotheses of Theorem~\ref{thm:ratpower}, fix
a degree \(e\ge1\). Without a height ceiling, the least
exponents on the height in normalized lower bounds of
the form \(e^{-\epsilon n}\) are
\begin{align*}
  \kappa_e &= e\frac{\log b}{\log a}
       &&\text{for nonzero polynomial values},\\
  \tau_e &= \max\!\left\{1,\ e+1-2\frac{\log(a/b)}{\log a}\right\}
       &&\text{for nonzero root distances}.
\end{align*}
\end{corollary}

\begin{proof}
Write \(A,B,t\) as in the preceding proof. For \(H(P)\le A\),
divide \eqref{eq:higher-ratpoly} by
\(H(P)^{e-\kappa_e}\le t^e\). For \(H(P)>A\), use
\(|P(t)|\ge B^{-e}\) and \(A^{\kappa_e}=B^e\).
This proves the polynomial bound.

If \(b>1\), B\'ezout's identity gives an integral polynomial
of degree at most \(e\) with \(B^eP(t)=1\), since \(A^e\)
and \(B^e\) are coprime. For each \(i=0,\ldots,e-1\)
in succession, add an integral multiple of \(Z^i(BZ-A)\)
to make the \(i\)-th coefficient at most \(A/2\) in
absolute value. Previous coefficients and \(P(t)\) are
preserved. The value at \(t\) then gives
\[
   |a_e|\le \frac{A}{2(t-1)}+B^{-e}t^{-e}.
\]
Thus \(H(P)\le A\) for large \(n\). The degree is \(e\),
since \(P(t)=B^{-e}\) in lowest terms, and the value
identity implies primitivity. When \(b=1\), instead take
\(P=Z^{e-1}(Z-A)+1\), of degree \(e\), with height at most
\(A\) and value \(1\). In both cases
\eqref{eq:higher-ratpoly} implies
\(\log H(P)=n\log a+o(n)\). The value is \(B^{-e}\),
so for every \(\kappa<\kappa_e\) these examples have
\(\log(|P(t)|H(P)^\kappa)/n\to(\kappa-\kappa_e)\log a<0\).

For the root bound, distances greater than \(1\) are
immediate since \(\tau_e\ge1\). If \(|t-\alpha|\le1\)
and \(H(\alpha)\le A\), use \eqref{eq:higher-ratroot}
and \(H(\alpha)^{e+1-\tau_e}\le t^2\).
If \(H(\alpha)>A\), rational evaluation of \(P_\alpha\)
and Lemma~\ref{lem:higher-transfer} give
\begin{align*}
 |t-\alpha|H(\alpha)^{\tau_e}
 &\gg_e t^{2-e}B^{-e}H(\alpha)^{\tau_e-1}\\
 &\ge t^{2-e}B^{-e}A^{\tau_e-1}
   = A^{\tau_e-e+1}B^{-2}\ge1 .
\end{align*}
A fixed constant is absorbed by the exponential tolerance.

For \(e=1\) and \(a\le b^2\), take a solution of
\(Aq-Bp=1\) with \(B\le q<2B\). Then \(p/q\) is reduced,
\(H(p/q)\asymp A\), and \(|t-p/q|\asymp B^{-2}\).
This shows sharpness of \(\tau_1=2\log b/\log a\) in this
case. If \(a>b^2\), then \(\tau_1=1\). For each \(n\),
the same congruence admits arbitrarily large \(q\), with
\(H(p/q)\asymp t q\) and \(|t-p/q|=1/(Bq)\).
Thus any exponent below \(1\) fails with unrestricted
heights.

For \(e\ge2\), fix a prime \(p\nmid ab\).
Start with a B\'ezout polynomial as above with \(B^eP(t)=1\).
In the coefficient reductions, choose the multiple of
\(Z^i(BZ-A)\) in a prescribed residue class modulo \(p^2\)
for \(i=0\) and modulo \(p\) for \(1\le i<e\).
This arranges \(a_0\equiv p\pmod{p^2}\) and
\(a_i\equiv0\pmod p\) for \(1\le i<e\), with
\(|a_i|\ll_p A\) for \(i<e\), since \(p\nmid A\).
The value identity gives \(a_e A^e\equiv1\pmod p\) and
\(|a_e|\ll_{e,p} A/t\). The polynomial is primitive
by the value identity and irreducible by Eisenstein's
criterion. Adding \(p^2Z^{e-1}(BZ-A)\) preserves all
these properties and changes the derivative at \(t\)
by \(p^2 A t^{e-2}\). One of the two choices therefore
has \(|P'(t)|\gg_p A t^{e-2}\).

On \(|z-t|\le1\), the coefficient bounds give
\(|P''(z)|\ll_{e,p} A t^{e-3}\). At
\(z=t-2P(t)/P'(t)\), the value predicted by the linear
Taylor terms is \(-P(t)\). The remainder divided by
\(P(t)\) is \(O_{e,p}(B^{-e}/(A t^{e-1}))=o(1)\).
Consequently the real interval between \(z\) and \(t\)
contains a root \(\alpha\) of degree \(e\) with
\[
 |t-\alpha|\ll_{e,p} A^{-(e-1)}B^{-2},\qquad
 H(\alpha)=H(P)\ll_{e,p} A .
\]
Since \(\tau_e=e-1+2\log b/\log a\) for \(e\ge2\),
any nonnegative smaller exponent gives an exponentially
small normalized distance. Negative exponents also fail
because \(H(\alpha)\ge1\) and the distances themselves
decay exponentially.
\end{proof}


\subsection{The cubic Pisot trace ratio}

Let \(h(Z)=Z^3-Z-1\) have roots \(\theta>1\) and
\(\beta,\gamma=\bar\beta\), and put
\[
 T_n=\theta^n+\beta^n+\gamma^n,\qquad R_n=T_{n+1}/T_n
\]
for all sufficiently large \(n\). The \(T_n\) are integers.
We use the splitting field and character lattice of
Proposition~\ref{prop:cubic-scalar}. Write \(K\) for the
splitting field, \(x_1=x,x_2=y,x_3=z\) for the characters
with values \(\theta^n,\beta^n,\gamma^n\), and
\(\lambda=\log\theta\). In particular \(xyz=1\), the
permutation group \(S_3\) acts on these characters, and
their lattice identifies with the centered integer
triples \((a_i-\bar a)_{i=1}^3\). For a Laurent polynomial,
\(w\) denotes the maximum of its monomial weights at the
embedding with dominant root \(\theta\), using the
ordinary complex modulus. Thus
\[
 w(x)=\lambda,\qquad w(y)=w(z)=-\lambda/2 .
\]
Use \(w(0)=-\infty\) and extend by subtraction to rational
functions. All character values are units. Let
\begin{gather*}
 \mathcal A=K[x^{\pm1},y^{\pm1},z^{\pm1}]/(xyz-1),\\
 V=x+y+z,\qquad U=\theta x+\beta y+\gamma z,\qquad R=U/V.
\end{gather*}
Brackets \([n]\) denote evaluation on this orbit. For each
of the three complex places \(v\) of \(K\), use the
product-formula normalized absolute value and monomial
weights \(w_v\). For an invariant function, \(w_v=w/3\)
at every such place. Nonzero Laurent polynomials and
rational functions have the local growth rates of
Lemma~\ref{lem:localweight}.

\begin{theorem}\label{thm:cubicquad}
For the cubic trace ratio and \(d=2\), the cutoffs in
\eqref{eq:higher-poly} and \eqref{eq:higher-root} are
both \(\theta\), with failure at equality.
With exact hits excluded, both estimates hold uniformly
without a height ceiling.
\end{theorem}

\begin{lemma}\label{lem:higher-cubic-local}
For every fixed \(C>0\), \eqref{eq:higher-poly} holds for
\(d=2\) and \(H(P)\le C\theta^{n/2}\).
\end{lemma}

\begin{proof}
The roots satisfy \(|\beta|=|\gamma|=\theta^{-1/2}\).
In particular \(T_n/\theta^n=1+O(\theta^{-3n/2})\) and
\begin{equation}\label{eq:higher-cubic-err}
 R_n-\theta=\frac{s_n}{T_n},\qquad
 s_n=(\beta-\theta)\beta^n+(\gamma-\theta)\gamma^n,\qquad
 R_n-\theta=O(\theta^{-3n/2}).
\end{equation}
For a nonzero polynomial \(P\) of degree at most two,
the algebraic integer \(P(\theta)\) has nonzero integral
norm. Since \(|P(\beta)|,|P(\gamma)|\ll H(P)\), this gives
\(|P(\theta)|\gg H(P)^{-2}\). Also
\(\lvert P(R_n)-P(\theta)\rvert\ll H(P)\theta^{-3n/2}\).
Consequently any \(\epsilon\)-violation of the polynomial
bound has \(H(P)\gg\theta^{n/2}\) and
\(\theta^n|P(R_n)|=O(e^{-\epsilon n})\).

Represent \(Z^nP\bmod h\) by an integral polynomial
\(B_n\) of degree at most two. It is nonzero since
\(Z\) is a unit in \(\mathbb Z[Z]/(h)\). At \(\beta,\gamma\)
its values are \(O(H(P)\theta^{-n/2})\). At \(\theta\) its
value is
\(O(\theta^n|P(R_n)|+H(P)\theta^{-n/2})\).
All three values are bounded under the height ceiling.
Interpolation in the fixed roots bounds the integral
coefficients. On any infinite sequence of violations
we can thus fix \(B=B_n\ne0\) on a subsequence.

For \(u\in\{\theta,\beta,\gamma\}\) let \(\ell_u\)
be the quadratic Lagrange polynomial taking value \(1\)
at \(u\) and \(0\) at the other two roots. Interpolating
\(P\) with values \(u^{-n}B(u)\) and using
\eqref{eq:higher-cubic-err} gives
\begin{align*}
 \theta^nP(R_n)
  = B(\theta)
  &+\ell_\beta'(\theta)B(\beta)\frac{s_n}{\beta^n}
   +\ell_\gamma'(\theta)B(\gamma)\frac{s_n}{\gamma^n}\\
  &+O_B(\theta^{-3n/2}).
\end{align*}
For the term at \(\theta\), use
\(\ell_\theta(R_n)=1+O(\theta^{-3n/2})\).
In the other terms, the second-order errors in
\(\ell_u(R_n)\) are \(O(\theta^{-3n})\), and the
multipliers \((\theta/u)^n\) have modulus
\(\theta^{3n/2}\). The factor \(\theta^n/T_n\)
in the first-order terms introduces an error of the
same size \(O_B(\theta^{-3n/2})\).

The terms before the error form a Laurent polynomial
in \((\beta/\gamma)^n\), with both nonconstant
coefficients nonzero. Indeed \(\ell_\beta'(\theta)\)
and \(\ell_\gamma'(\theta)\) are nonzero because their
zeros at \(\theta\) are simple, and \(B\) cannot vanish
at a root of the irreducible cubic \(h\). The ratio
\(\beta/\gamma\) is not a root of unity by the character
lattice, and has modulus \(1\). Lemma~\ref{lem:localweight}
gives logarithmic size \(o(n)\) for this nonzero Laurent
polynomial, contrary to the exponentially small
\(\theta^nP(R_n)\) and error term. This proves the lemma.
\end{proof}

\begin{proof}[Proof of Theorem~\ref{thm:cubicquad}]
Suppose there are infinitely many violations of
\eqref{eq:higher-poly} for some \(Q<\theta\) and a fixed
positive tolerance. Divide each \(P\) by its content,
which only decreases the normalized absolute value
and the height. On a subsequence, with \(s\in[0,\lambda)\)
and \(\delta>0\), the primitive coefficient vectors
\({\bf a}_n=(a_0,a_1,a_2)\) satisfy
\begin{equation}\label{eq:higher-cubic-size}
 \log H(P)/n\ \longrightarrow\ s,\qquad
 |P(R_n)|\le \exp((-2s-\delta)n).
\end{equation}
We first show how to obtain formal relations among
these vectors and how to descend inside a plane.

Write \(F_K=\operatorname{Frac}(\mathcal A)\).
Use either the whole space \(F_K^3\) or a fixed plane
of relation \(A\cdot{\bf a}=0\), assumed to contain the
vectors after evaluation on the subsequence. In the
plane case assume the relation normal
\(A=(A_0,A_1,A_2)\) is a primitive, componentwise invariant
Laurent tuple. Here primitive means that its nonzero
components have no common nonunit Laurent factor.
Define for the plane
\begin{equation}\label{eq:higher-cubic-k}
 k_A=\max\!\left\{
    w(A_0)-2s-\delta,\ 
    s+\max_{i=1,2} w(A_i-R^i A_0)
 \right\}.
\end{equation}
We claim that a new formal relation holds on a further
infinite subsequence in the whole-space case, and also in
the plane case if \(k_A>-\delta\). In the plane case
the new relation is independent of \(A\) over \(F_K\).

To prove the claim, use the forms \(X^\mu a_i\),
\(0\le i\le2\), for \(\mu\in[0,H]^2\) in a fixed
basis of the character lattice. The notation \(X^\mu\)
here denotes a character, and \(M=(H+1)^2\) counts
the shifts. Put \(W_i=U^iV^{2-i}\) for \(i=0,1,2\).
At a complex place \(v\), take the leading monomial
\(X^{\mu_0}\) of \(W_0=V^2\) by \(w_v\), breaking
ties by a translation-invariant lexicographic order.
Replace each form \(X^\mu a_0\) for which the shifted
supports fit inside the box by
\[
                  X^{\mu-\mu_0}\sum_i W_i a_i .
\]
There are \(M+O(H)\) such slots. The change of basis
is triangular on block \(0\) with nonzero pivots,
and does not change blocks \(1\) and \(2\).
Assign weights \(w_v(X^\mu)+(-2s-\delta)/3\) to
the replacements, \(w_v(X^\mu)+s/3\) to the retained
forms, and use the maximum of the weights of the
nonzero basis terms as the weight of a form.
The weights bound the logarithms of evaluations
divided by \(n\), up to \(o(1)\), by
\eqref{eq:higher-cubic-size} and
\(w_v(X^{-\mu_0}V^2)=0\). This also bounds the
evaluation of any fixed linear combination of
these forms.

The old basis with weights \(w_v(X^\mu)\) and the
new weighted basis give weights differing by \(O(1)\)
on every nonzero form, uniformly in \(H\). In the
forward expansion all support shifts are fixed.
For inverse expansion the substitutions remain
triangular on block \(0\), moving only to smaller
monomials there, with weight never increasing
within that block. Each such step introduces
terms in the other blocks with a fixed bound on
their weight increase. Those blocks require no
further substitution. The fixed additions to
basis weights therefore preserve the uniform bound.

In the plane case the relation space inside the box
is exactly the span of \(X^\nu A\) for \(\nu\) in a
rectangle of \(M+O(H)\) slots with fixed margins.
Indeed a Laurent tuple in the \(F_K\)-span of \(A\)
is a Laurent multiple of \(A\), by primitivity and
unique factorization. The minimum and maximum
exponents of each individual character coordinate
add under multiplication of Laurent polynomials.
Applying this in both directions of both axes gives
the rectangle condition on all monomials of the
multiplier. In a rectangle shrunk by another fixed
margin, consider the relation forms
\(X^{\nu-\mu_0}W_0 A\). The margin ensures their
supports fit, and all shifted slots needed for
replacing the block-\(0\) terms fit as well. These
forms span a subspace of codimension \(O(H)\) in
the relation space. Under the old weights their
leading pivots are distinct translates, with
weights \(w_v(X^\nu)+w_v(A)\). Under the new basis
the block-\(0\) factors are the coefficients of
\(X^\nu A_0\) on the replacement slots; the other
blocks have factors
\[
              X^{\nu-\mu_0}(W_0 A_i-W_i A_0).
\]
Their leading pivots translate with \(\nu\).
Since each \(A_i-R^i A_0\) is invariant and
\(w_v(X^{-\mu_0}W_0)=0\), the new weights are
\(w_v(X^\nu)+k_A/3\).
Distinct pivots in both calculations follow by
ordering basis terms first by weight and then
by a translation-invariant order on monomials
and block indices.

For clarity, the increasing list of weights of
a filtered space is given by the jumps of the
dimension of its weight sublevels, counted with
multiplicity. A subspace receives the induced
filtration, and its quotient receives the
filtration by minimum weight among lifts.
The lists for a subspace and the quotient
together give the list for the whole space,
by additivity of dimensions at each sublevel.
Two weights at uniformly bounded distance
give ordered lists at bounded distance termwise,
also on an induced quotient. Consequently the
\(O(H)\) boundary dimensions of the relation space
contribute only \(O(H)\) to the change of
its weight sum. The interior forms give the
change computed above. Before replacement, the
weight sum of the relation space is the sum of
the shifted weights \(w_v(X^\nu)+w_v(A)\) on the
full rectangle, since the shifts have distinct
leading pivots. Summing over the complex places
cancels all monomial weights by the product
formula and the unit property. It follows that
the sum of the new relation weights over those
places is \(M k_A+O(H)\).

The ambient weight sum over the same places
is \(-M\delta+O(H)\), again by cancellation of
the monomial weights. Choose a basis in the
quotient at each place by taking classes of
ambient basis forms in order of increasing
weight and skipping dependent classes.
Its total weight is the ambient sum minus the
relation sum. In the whole-space case the
relation sum is zero. In the plane case the
total quotient weight over the places is
\(M(-\delta-k_A)+O(H)\). In either case of the
claim it is negative for large fixed \(H\).
Use a fixed basis of the quotient selected
from the \emph{original} monomial-coordinate
classes to form the evaluation vector.
It has algebraic integer entries, is nonzero
since \({\bf a}_n\ne0\), and has height at most
exponential in \(n\). The chosen local bases
give a product exponentially small in \(n\).
The number-field Subspace Theorem
\cite{Schlickewei} gives a fixed nonzero
linear equation on these quotient coordinates
on a further infinite subsequence. As a formal
relation it is outside the previous relation
space by the choice of a quotient basis.
This proves the claim.

All relations can be intersected with their
Galois conjugates because the coefficient vectors
are rational and evaluation commutes with the
Galois action. The resulting coefficient spaces
are invariant, and the ranks of the resulting
formal systems stabilize on evaluation, by
Lemma~\ref{lem:localweight} on nonzero minors.
After one relation we therefore have a line
or a plane, since the zero space cannot contain
these vectors. For an invariant line or relation
normal, a primitive Laurent representative
is componentwise invariant after a common
unit rescaling. To see that the normalization
in Proposition~\ref{prop:cubic-scalar} applies to
tuples, take the union of the supports for its
translation step. Primitivity makes the Galois
multipliers common Laurent units, and the
barycenter of the union gives the same lattice
translation as for a pair. The remaining common
scalar cocycle is removed by Hilbert's Theorem~90.
In particular the maximum weight of a nonzero
invariant tuple is a nonnegative integer multiple
of \(\lambda/2\).

Suppose the vectors lie on an invariant line
with such a representative \(A\), and let
\(D=\max_i w(A_i)\). Its primitive integral
specialization has logarithmic height
\(Dn+o(n)\). Here is the gcd justification.
Clear denominators of the coefficients by a fixed
integer, so the evaluated coordinates are rational
algebraic integers. A nonzero coordinate and a
fixed linear combination avoiding its irreducible
factors form a coprime Laurent pair. After a common
monomial shift the pair consists of coprime
polynomials. The shift is a unit on the orbit
at all finite places. Off a fixed finite set of
places the torus gcd theorem
\cite[Theorem~1.13(a)]{LW} gives
logarithmic gcd \(o(n)\), as the orbit meets each
exceptional coset finitely often by its torsion-free
character lattice. At the finitely many remaining
primes every character is a unit, so
Lemma~\ref{lem:localweight} gives logarithmic gcd
\(o(n)\) there too. At the chosen embedding the
maximum coordinate has logarithmic size \(Dn+o(n)\)
by the same local-weight lemma. This proves the
height assertion. Since the vectors are primitive,
\eqref{eq:higher-cubic-size} now gives \(D=s<\lambda\)
and hence \(D\le\lambda/2\). The fixed Laurent
representative has coordinates \(O(\theta^{n/2})\)
at evaluation by the triangle inequality, so
the primitive vectors also have
\(H(P)\ll\theta^{n/2}\), contrary to
Lemma~\ref{lem:higher-cubic-local}.

It remains to show that every invariant plane
produced above satisfies \(k_A>-\delta\).
Suppose instead \(k_A\le-\delta\). Then \(A_0\ne0\):
otherwise each nonzero \(A_i\) would have
\(w(A_i)\le-s-\delta<0\), which is impossible
for an invariant Laurent polynomial. Since
\(w(R)=0\) and \(w(A_0)\ge0\), the inequalities
in \eqref{eq:higher-cubic-k} give
\begin{equation}\label{eq:higher-cubic-support}
 D=\max_i w(A_i)=w(A_0)\le2s<2\lambda,\qquad
 w(A_m-R^m A_0)\le-s-\delta\quad(m=1,2).
\end{equation}
Thus \(D\le3\lambda/2\). A centered exponent triple
\({\bf c}\) has ordinary weight \(3\lambda c_1/2\).
Since every permutation is in the joint support
whenever \({\bf c}\) is, all its coordinates are
at most \(1\). The only allowed orbits are
\[
  1,\quad x_i,\quad x_i^{-1},\quad x_i^{-2},\quad
  x_i^{-3},\quad x_i/x_j\ (i\ne j).
\]
Indeed the coordinates of a centered triple sum
to zero and have common fractional part \(0\),
\(1/3\), or \(2/3\). For fractional part \(0\)
the possible triples up to permutation are
\((0,0,0),(1,0,-1),(1,1,-2)\).
For fractional part \(1/3\) the only possibility
is \((1/3,1/3,-2/3)\).
For fractional part \(2/3\), the possibilities
are \((2/3,-1/3,-1/3)\) and
\((2/3,2/3,-4/3)\).

Use the first-order expansion at the chosen embedding
\begin{equation}\label{eq:higher-cubic-expansion}
 R^m=\theta^m+
  m\theta^{m-1}\bigl((\beta-\theta)y/x+(\gamma-\theta)z/x\bigr)
  +O_w(-3\lambda)\qquad(m=1,2),
\end{equation}
where \(O_w(u)\) denotes a rational function of
weight at most \(u\). This follows by using
\(w(y/x)=w(z/x)=-3\lambda/2\) in
\(R=(\theta+\beta y/x+\gamma z/x)/(1+y/x+z/x)\).
After multiplication by \(A_0\), the remainder
has weight at most \(-3\lambda/2\). Thus every
comparison of Laurent coefficients below is made
at weights exceeding the remainder.

All positive-weight coefficients of \(A_m\) must
agree with those of \(\theta^m A_0\), because the
first-order correction times \(A_0\) has weight
at most zero. In particular this applies to \(y^{-j}\)
for \(1\le j\le3\). Each coefficient there is fixed
by the stabilizer of \(\beta\), so lies in
\(\mathbb Q(\beta)\), which does not contain \(\theta\).
Indeed equality of the two cubic fields would
make \(\mathbb Q(\beta)=\mathbb Q(\theta)\) real.
The comparison with \(m=1\) forces the \(A_0\)
coefficient to vanish, and both comparisons
force the \(A_m\) coefficients to vanish. Galois
invariance removes all three inverse-root orbits.

Write \(c_{ij}\) for the coefficient of \(x_i/x_j\)
in \(A_0\). At \(x/y,x/z\) the positive-weight
comparison gives coefficients \(\theta^m c_{1j}\)
in \(A_m\). Conjugating gives coefficients
\(\alpha_i^m c_{ij}\) throughout the ratio orbit,
where \((\alpha_1,\alpha_2,\alpha_3)=(\theta,\beta,\gamma)\).
At \(y/z\), of weight zero, comparison in
\eqref{eq:higher-cubic-expansion} gives
\[
 (\beta^m-\theta^m)c_{23}
 -m\theta^{m-1}(\beta-\theta)c_{13}=0
 \qquad(m=1,2).
\]
The term from \(z/x\) would require \(xy/z^2=z^{-3}\)
in \(A_0\), which was already excluded. For \(m=1\)
we get \(c_{23}=c_{13}\), and then \(m=2\) gives
\((\beta-\theta)^2c_{23}=0\). Invariance removes
the whole ratio orbit.

Only \(1,x,y,z\) can now appear, so \(D\le\lambda\)
and the first-order corrections times \(A_0\)
have strictly negative weight. Comparison at
weight zero forces the constant coefficients
to vanish, since they are rational and
\(\theta\) is irrational. A nonzero direct-root
orbit would have \(D=\lambda\) and hence
\(s\ge\lambda/2\). Comparing at \(x\) and then
conjugating gives coefficients \(\alpha_i^m c_i\)
in \(A_m\), where \(c_i\) is the coefficient
of \(x_i\) in \(A_0\). At \(y\), of weight
\(-\lambda/2>-s-\delta\), comparison gives
\[
 (\beta^m-\theta^m)c_2
 -m\theta^{m-1}(\beta-\theta)c_1=0
 \qquad(m=1,2).
\]
The term from \(z/x\) would require \(xy/z=z^{-2}\),
already excluded. These two equations force
\(c_1=c_2=0\), which rules out a nonzero direct-root
orbit. We have contradicted \(A_0\ne0\).
Thus the plane descends to a line, already
excluded, and the polynomial bound holds at
every \(Q<\theta\). The root bound follows by
Lemma~\ref{lem:higher-transfer} since \(R_n\)
is bounded.

At the endpoint, exact hits occur infinitely
often within the height ceiling. The recurrence
\(T_{n+3}=T_{n+1}+T_n\) starts with \(3,0,2\).
The triples modulo \(2\), starting at \(n=0\),
cycle through
\[
 100,\ 001,\ 010,\ 101,\ 011,\ 111,\ 110 .
\]
In particular \(T_n,T_{n+1}\) are both even when
\(n\equiv1\pmod7\). For large such \(n\) the
reduced height of \(R_n\) is at most
\(T_{n+1}/2=(\theta/2+o(1))\theta^n<\theta^n\),
where \(\theta<2\). The primitive linear
polynomial of \(R_n\) has the same height.
These exact hits also give failure at every \(Q>\theta\).

Finally, if \(P(R_n)\ne0\) and \(\deg P\le2\),
rational evaluation gives
\[
 |P(R_n)|\ge \operatorname{den}(R_n)^{-2}
             \ge |T_n|^{-2},
\]
where \(\operatorname{den}\) denotes the reduced
positive denominator. Any nonexact violation
with tolerance \(\epsilon>0\) therefore has
\(H(P)<|T_n|e^{-\epsilon n/2}\).
This fits a fixed ceiling \(Q^n\) with \(1<Q<\theta\)
for large \(n\), and is impossible by the bound
just proved. Lemma~\ref{lem:higher-transfer} gives
the unrestricted root bound for \(\alpha\ne R_n\).
\end{proof}

\begin{proposition}[Height-one failures]\label{prop:higher-height-one}
For \(R_n=F_{n+1}/F_n\), where \(F_n\) are the Fibonacci
numbers, both bounds fail at \(Q=1\) for every \(d\ge2\).
For the cubic trace ratio above, both fail at \(Q=1\)
for every \(d\ge3\). The failures persist with exact
hits excluded.
\end{proposition}

\begin{proof}
For the Fibonacci ratio let
\(\varphi=(1+\sqrt5)/2\) and \(\psi=-1/\varphi\).
Binet's formula gives \(F_{n+1}-\varphi F_n=\psi^n\),
so \(0<|R_n-\varphi|\ll\varphi^{-2n}\). The number
\(\varphi\) and its minimal polynomial \(Z^2-Z-1\)
both have height \(1\), and the polynomial value
at \(R_n\) is \(O(\varphi^{-2n})\) by factorization.
For the cubic ratio use \(\theta\) and \(h=Z^3-Z-1\),
also both of height \(1\). Equation
\eqref{eq:higher-cubic-err} and boundedness of \(R_n\)
give \(|R_n-\theta|,|h(R_n)|\ll\theta^{-3n/2}\).
The rational targets cannot equal these irrational
approximants or be zeros of their minimal polynomials.
In each case take any positive tolerance smaller
than the rate of exponential decay.
\end{proof}

\begin{thebibliography}{99}
\bibitem{BCZ} Y.~Bugeaud, P.~Corvaja and U.~Zannier, \href{https://doi.org/10.1007/s00209-002-0449-z}{An upper bound for the G.C.D.\ of \(a^n-1\) and \(b^n-1\)}, Math. Z. \textbf{243} (2003), 79--84.

\bibitem{BE} Y.~Bugeaud and J.-H.~Evertse, \href{https://www.numdam.org/item/ASNSP_2009_5_8_2_333_0/}{Approximation of complex algebraic numbers by algebraic numbers of bounded degree}, Ann. Sc. Norm. Super. Pisa Cl. Sci. (5) \textbf{8} (2009), 333--368.

\bibitem{CZ04} P.~Corvaja and U.~Zannier, \href{https://doi.org/10.1007/BF02392563}{On the rational approximations to the powers of an algebraic number: solution of two problems of Mahler and Mend\`es France}, Acta Math. \textbf{193} (2004), 175--191.

\bibitem{CZgcd} P.~Corvaja and U.~Zannier, \href{https://arxiv.org/abs/math/0311030}{A lower bound for the height of a rational function at \(S\)-unit points}, Monatsh. Math. \textbf{144} (2005), 203--224.

\bibitem{CZ05} P.~Corvaja and U.~Zannier, \href{https://doi.org/10.5802/jtnb.517}{On the length of the continued fraction for values of quotients of power sums}, J. Th\'eor. Nombres Bordeaux \textbf{17} (2005), 737--748.

\bibitem{CZ08} P.~Corvaja and U.~Zannier, \href{https://doi.org/10.1007/978-3-211-74280-8_8}{Applications of the Subspace Theorem to certain Diophantine problems: a survey of some recent results}, in \emph{Diophantine Approximation: Festschrift for Wolfgang Schmidt}, Dev. Math. \textbf{16}, Springer, Vienna, 2008, 161--174.

\bibitem{CZ18} P.~Corvaja and U.~Zannier, \href{https://doi.org/10.1017/9781108348096}{\emph{Applications of Diophantine Approximation to Integral Points and Transcendence}}, Cambridge Tracts in Mathematics \textbf{212}, Cambridge University Press, 2018.

\bibitem{EF} J.-H.~Evertse and R.~G.~Ferretti, \href{https://annals.math.princeton.edu/2013/177-2/p04}{A further improvement of the quantitative Subspace Theorem}, Ann. of Math. (2) \textbf{177} (2013), 513--590.

\bibitem{FH} C.~Fuchs and S.~Heintze, \href{https://doi.org/10.5802/jtnb.1247}{Approximation of values of algebraic elements over the ring of power sums}, J. Th\'eor. Nombres Bordeaux \textbf{35} (2023), 63--86.

\bibitem{HRI} Harish-Chandra Research Institute, \href{https://www.hri.res.in/~anlweb/Academic-Report-2024-25.pdf}{\emph{Academic Report 2024--25}}, research summary of G.~Prasad, p.~59.

\bibitem{ambient} Z.~Hu and B.~Li, \href{https://doi.org/10.1007/s10473-025-0422-1}{Hausdorff dimension of recurrence sets for matrix transformations of tori}, Acta Math. Sci. Ser. B \textbf{45} (2025), 1659--1673. \href{https://arxiv.org/abs/2402.04810}{Author preprint}.

\bibitem{KL} R.~Kannan and L.~Lov\'asz, \href{https://annals.math.princeton.edu/1988/128-3/p06}{Covering minima and lattice-point-free convex bodies}, Ann. of Math. (2) \textbf{128} (1988), 577--602.

\bibitem{KP} V.~Kumar and G.~Prasad, \href{https://doi.org/10.1515/forum-2024-0441}{On the rational approximation to linear combinations of powers}, Forum Math. \textbf{38} (2026), 1027--1041.

\bibitem{KPsub} V.~Kumar and G.~Prasad Sena, \emph{Thue--Roth's type inequality beyond the \(S\)-units}, revised manuscript submitted to Ann. Sc. Norm. Super. Pisa Cl. Sci., listed on \href{https://sites.google.com/view/veekeshkumar/publications}{V.~Kumar's publications page}; see also \cite{HRI}.

\bibitem{KT} V.~Kumar and R.~Thangadurai, \href{https://doi.org/10.1112/mtk.12161}{On simultaneous approximation of algebraic numbers}, Mathematika \textbf{68} (2022), 1153--1175.

\bibitem{LechSML} C. Lech,
\href{https://doi.org/10.1007/BF02590997}{A note on recurring series},
Ark. Mat. \textbf{2} (1953), no.~5, 417--421.

\bibitem{Levin} A.~Levin, \href{https://doi.org/10.1007/s00222-018-0831-z}{Greatest common divisors and Vojta's conjecture for blowups of algebraic tori}, Invent. Math. \textbf{215} (2019), 493--533.

\bibitem{LW} A.~Levin and J.~T.-Y.~Wang, \href{https://doi.org/10.1515/crelle-2019-0033}{Greatest common divisors of analytic functions and Nevanlinna theory on algebraic tori}, J. Reine Angew. Math. \textbf{767} (2020), 77--107, published online in 2019.

\bibitem{Hilbert90Milne} J. S. Milne,
\href{https://www.jmilne.org/math/CourseNotes/CFT.pdf}{Class Field Theory},
course notes, version~4.03, 2020.

\bibitem{Northcott1950} D. G. Northcott,
\href{https://doi.org/10.2307/1969504}{Periodic points on an algebraic variety},
Ann. of Math. (2) \textbf{51} (1950), no.~1, 167--177.

\bibitem{RV} M.~Ru and P.~Vojta, \emph{Schmidt's Subspace Theorem with moving targets}, Invent. Math. \textbf{127} (1997), 51--65.

\bibitem{Schlickewei} H.~P.~Schlickewei, \href{https://www.numdam.org/item/CM_1992__82_3_245_0/}{The quantitative subspace theorem for number fields}, Compositio Math. \textbf{82} (1992), 245--273.

\bibitem{Schmidt70} W.~M.~Schmidt, \href{https://doi.org/10.1007/BF02392334}{Simultaneous approximation to algebraic numbers by rationals}, Acta Math. \textbf{125} (1970), 189--201.

\bibitem{Wirsing} E.~Wirsing, \href{https://doi.org/10.1515/crll.1961.206.67}{Approximation mit algebraischen Zahlen beschr\"ankten Grades}, J. Reine Angew. Math. \textbf{206} (1961), 67--77.

\bibitem{Xiao} Z.~Xiao, \href{https://doi.org/10.1007/s00209-024-03453-4}{Greatest common divisors for polynomials in almost units and applications to linear recurrence sequences}, Math. Z. \textbf{306} (2024), Paper No.~61. \href{https://arxiv.org/abs/2110.01751v3}{Author's preprint}.

\end{thebibliography}
\end{document}
